\documentclass[11pt,reqno]{amsart}

%==============================================================
% Primitive Divergence and Observable It\^o Calculus over
% Nonclassical Noises
% --------------------------------------------------------------
% Companion to the IBM trilogy:
%   Fontes, Energy-First Stochastic Calculus       [EnergyFirst]
%   Fontes, First-Chaos Dual Closure of the IBM... [FCDC]
%   Fontes, A Bidegree Atlas for the IBM ...       [BidegreeAtlas]
% Scope: framework paper on the operator-derivative principle over
% nonclassical noise. No new black-noise theorems are proved; no
% primitive divergence is constructed.
%==============================================================

\usepackage[utf8]{inputenc}
\usepackage[T1]{fontenc}
\usepackage{lmodern}
\usepackage[margin=1.15in]{geometry}
\usepackage{amsmath,amssymb,amsthm,mathtools}
\usepackage{mathrsfs}
\usepackage{microtype}
\usepackage{hyperref}
\usepackage{cleveref}

\emergencystretch=3em

% --- theorem environments (compatible with the trilogy)
\theoremstyle{plain}
\newtheorem{theorem}{Theorem}[section]
\newtheorem{proposition}[theorem]{Proposition}
\newtheorem{lemma}[theorem]{Lemma}
\newtheorem{corollary}[theorem]{Corollary}

\theoremstyle{definition}
\newtheorem{definition}[theorem]{Definition}
\newtheorem{example}[theorem]{Example}
\newtheorem{assumption}[theorem]{Assumption}
\newtheorem{program}[theorem]{Program}

\theoremstyle{remark}
\newtheorem{remark}[theorem]{Remark}
\newtheorem*{sketch}{Sketch}

% --- macros
\newcommand{\R}{\mathbb{R}}
\newcommand{\N}{\mathbb{N}}
\newcommand{\E}{\mathbb{E}}
\newcommand{\Prob}{\mathbb{P}}
\newcommand{\Dom}{\mathrm{Dom}}
\newcommand{\cl}{\mathrm{cl}}
\newcommand{\IBM}{Z}

\title{Primitive Divergence and Carr\'e-du-Champ Seeds\\
       over Nonclassical Noises}
\author{Ramiro Fontes}
\date{\today}

\begin{document}

\begin{abstract}
This paper formulates the operator-derivative problem over
nonclassical noises. The primitive datum is a tuple
$(\{\Gamma_{\mathcal N,S}\}_{|S|\ge 2},\delta_{\mathcal N},D_{\mathcal N})$:
a generator-level interaction hierarchy
$\{\Gamma_{\mathcal N,S}\}_{|S|\ge 2}$ of $|S|$-linear
product-defect forms indexed by finite subsets $S$, on a test
algebra over a Tsirelson--Vershik noise $\mathcal N$, headed by the
bilinear two-point carr\'e-du-champ
$\Gamma_{\mathcal N}^{(2)}=\Gamma_{\mathcal N,\{i,j\}}$;
a closed primitive divergence $\delta_{\mathcal N}$ on
$L^{2}(\mathcal N)$; and the adjoint operator derivative
$D_{\mathcal N}:=\delta_{\mathcal N}^{*}$.
Classical stochastic calculus starts with a first-chaos innovation
and derives $\delta$ from it; the operator-derivative framework
reverses this, treating
$(\{\Gamma_{\mathcal N,S}\}_{|S|\ge 2},\delta_{\mathcal N})$
as primitive data and Brownian, Poisson, and L\'evy stochastic
integrals as realizations of $\delta$ rather than its definition.
We make four contributions.
\textbf{(1)} \emph{Representation-free primitive.} The framework is
set up over $\mathcal N$ without reference to a first-chaos
generator, so that classical innovations are particular realizations
rather than definitional inputs, and the candidate primitive datum
is the generator-level interaction hierarchy
$\{\Gamma_{\mathcal N,S}\}_{|S|\ge 2}$, headed by the bilinear,
positive two-point carr\'e-du-champ
$\Gamma_{\mathcal N}^{(2)}=\Gamma_{\mathcal N,\{i,j\}}$, together
with the closed divergence $\delta_{\mathcal N}$.
\textbf{(2)} \emph{Seed-collapse theorem.} Every centered
interval-additive deterministic seed of a candidate
$\delta_{\mathcal N}$ lies in the first chaos of $\mathcal N$, and
hence is trivial over a black noise; the positive replacement is a
carr\'e-du-champ / product-defect seed.
\textbf{(3)} \emph{Leibniz-defect invariant.} The bilinear form
$\mathfrak L_D(F,G)=D(FG)-F\,DG-G\,DF$ is identified as the first
algebraic invariant of a nonclassical operator derivative, with
Wiener-zero and compensated-Poisson-quadratic benchmarks; over a
black noise the defect question becomes structural.
\textbf{(4)} \emph{Nonclassical divergence problem.} The
construction, or impossibility, of the primitive tuple
$(\{\Gamma_{\mathcal N,S}\}_{|S|\ge 2},\delta_{\mathcal N},D_{\mathcal N})$
intrinsically from the algebra of $\mathcal N$, without first-chaos
generators, is formulated as the central open problem; the two-point
form $\Gamma_{\mathcal N}^{(2)}$ is the primary candidate seed, and
$\mathfrak L_{D_{\mathcal N}}$ is the algebraic remainder relating
$\Gamma_{\mathcal N}^{(2)}$ to the candidate generator
$L_{\mathcal N}=-\delta_{\mathcal N}D_{\mathcal N}$, when this
composition is defined on a suitable dense domain. The classical
no-innovation-reduction statement appears as a short benchmark
corollary. No primitive divergence is constructed, no universal
existence theorem for the hierarchy
$\{\Gamma_{\mathcal N,S}\}_{|S|\ge 2}$ over black noises is
claimed, no representation of the hierarchy through
$(\delta_{\mathcal N},D_{\mathcal N})$ is given, and no new
black-noise theorem beyond the seed-collapse statement is asserted.
\end{abstract}

\maketitle

%==============================================================
\section{Introduction}\label{sec:intro}
%==============================================================

An \emph{observable It\^o identity} for a process $X$ on a Polish
space $E$ is a martingale-problem relation of the form
\[
  f(X_t) - f(X_0) = \int_0^t Lf(X_s)\,ds + M_t[f], \qquad f\in\mathcal T,
\]
where $L\colon\mathcal T\to C_b(E)$ is linear and $(M_t[f])$ is a
martingale with respect to the natural filtration of $X$. When the
martingale problem $(L,\mathcal T)$ is well posed, such an identity
fixes the semigroup, generator, and law of $X$. It does not, by
itself, fix how the randomness driving $M[f]$ is organized: the same
semigroup is compatible with many different choices of underlying
noise, of pathwise lift, and of operator-valued representation. The
purpose of this paper is to formulate the operator-derivative
problem over nonclassical noise and to identify its first two
structural invariants: the seed data of a candidate primitive
divergence, and the Leibniz defect of its adjoint.

The trilogy that motivates this paper ---
\cite{EnergyFirst,FCDC,BidegreeAtlas} --- adopts what we call the
\emph{energy-first stance}: the divergence operator $\delta$, together
with its adjoint $D=\delta^{*}$, is taken as primitive. Stochastic
integration against a Brownian or Poisson generator is then one
realization of $\delta$, not its definition. The trilogy uses this
stance to handle Iterated Brownian Motion (IBM), where the natural
$L^2$ closure of the pathwise It\^o correction fails on a first-chaos
sector and only the $H^{-1}$ dual closure survives. The framework was
developed precisely because the classical It\^o integral was not the
right primitive for that problem.

\subsection{Black noise as a test of the energy-first stance}\label{sec:intro-blacknoise}

The IBM obstructions of the trilogy live entirely inside a
\emph{classical} noise: the source randomness is a pair of
independent Brownian motions, and the obstruction is one of
topology and of sector summability, not of the noise algebra itself.
Tsirelson's theory of noises~\cite{Tsirelson2004,TsirelsonVershik1998}
exhibits a strictly more severe phenomenon. A \emph{noise} is a
continuous product of probability spaces, encoded as a family
$\mathcal N=(\mathcal F_{s,t})_{s\le t}$ of $\sigma$-fields with
$\mathcal F_{r,t}=\mathcal F_{r,s}\vee\mathcal F_{s,t}$ under
independence, and an associated time-shift. Inside any noise, the
\emph{classical subnoise} $\mathcal N^{\mathrm{cl}}$ is the noise
generated by additive integrable families $(X_{s,t})$ with
$X_{r,t}=X_{r,s}+X_{s,t}$; equivalently, by the closed linear span of
the first chaoses $I_1(\mathcal F_{s,t})$. The noise $\mathcal N$ is
\emph{classical} when $\mathcal N^{\mathrm{cl}}=\mathcal N$ and
\emph{black} when $\mathcal N^{\mathrm{cl}}$ is trivial while
$\mathcal N$ itself is not. Known examples of black noise include
Tsirelson's sticky and coalescing flows~\cite{LeJanRaimond2004}, the
noise of the Brownian web in the sense of~\cite{Tsirelson2004,SSS2017},
and the noise of two-dimensional critical
percolation~\cite{SchrammSmirnov2011}.

The role of black noise in the present paper is precisely calibrated.
It is straightforward, and we record it as a proposition in
\cref{sec:proposition}, that black noise forbids any reduction of the
martingale family $(M_t[f])$ to a stochastic integral against a
Brownian or Poisson driver living inside the filtration of $X$: there
is simply no classical first-chaos generator in which such an
integral could be written. What does \emph{not} follow from this, and
what is sometimes implicitly assumed in the literature, is that black
noise obstructs an energy-first It\^o calculus as such. Under the
energy-first stance, $\delta$ is primitive; the absence of a
classical innovation eliminates one realization of $\delta$ but does
not, prima facie, eliminate the operator. Whether $\delta$ admits a
genuinely intrinsic construction over a nonclassical noise --- without
an intermediate classical generator --- is the central programmatic
question of this paper, and we leave it open.

\subsection{Three calibrations}\label{sec:intro-calibrations}

The relative position of the trilogy and of the black-noise setting
is clearest when read against the Brownian benchmark.

\emph{Brownian (Wiener) noise.} The first chaos
$I_1(\mathcal F_{s,t})\simeq L^2([s,t])$ is the canonical home of the
additive integrable families; the Brownian increment
$W_t-W_r=(W_s-W_r)+(W_t-W_s)$ is the prototypical such family. The
classical It\^o divergence
$\delta^W(\mathbf 1_{[s,t]})=W_t-W_s$ is the seed of all subsequent
construction: by closure and chaos decomposition it extends from
$I_1$ to all of $L^2(\mathcal F_W)$, and the Malliavin pair
$(\delta^W,D^W)$ is determined accordingly.

\emph{IBM and the trilogy.} The natural noise of IBM is a tensor
product of two independent Brownian noises; it is classical, and the
first chaos is rich. The trilogy's obstructions arise \emph{inside}
this classical setting. In~\cite{EnergyFirst} the $L^2$ closure of
the mollified pathwise correction fails on a first-chaos sector and
only an $H^{-1}$ dual closure survives; in~\cite{FCDC} that dual
closure is identified explicitly; in~\cite{BidegreeAtlas} the higher
bidegree sectors $(m,n)$ are organized into an atlas with their own
summability conditions. None of these obstructions is
filtration-level: the seed of the classical construction is intact,
and the failure occurs in the closure and the cross-sector
summability.

\emph{Black noise.} Here the first chaos itself collapses:
$I_1(\mathcal N)$ is trivial, modulo constants. The Brownian template
breaks at the first step. There is no $\mathbf 1_{[s,t]}$ to integrate
against a classical generator, no isometry
$L^2([s,t])\to I_1(\mathcal F_{s,t})$ to extend by closure, and no
chaos decomposition into $\bigoplus_n I_n$ with $I_n$ built from
tensor powers of $I_1$. Any candidate energy-first construction must
therefore originate elsewhere, at the algebra of the noise rather than
at a first-chaos generator.

Thus IBM and black noise are first-class examples of the
operator-derivative framework along different axes: IBM tests the
topology of the correction inside classical noise, while black noise
tests whether the primitive divergence $\delta_{\mathcal N}$ can be
specified when the classical first-chaos seed is absent. The two
tests are orthogonal, and their resolutions are independent.

\subsection{The first structural pressure point}\label{sec:intro-seed}

The collapse just described is not by itself a theorem of
obstruction. What we record in \cref{sec:seed} is the precise
location of the pressure point. Let $\delta_{\mathcal N}$ be assembled,
in the manner of the classical case, from a family of interval
operators
\[
  \delta_{s,t}\colon\Dom(\delta_{s,t})\subset L^2(\mathcal F_{s,t};H_{s,t})\to L^2(\mathcal F_{s,t}),
\]
with test-direction spaces forming an additive family
$H_{r,t}=H_{r,s}\oplus H_{s,t}$. We call $\delta_{\mathcal N}$
\emph{interval-compatible} if its deterministic-coefficient seed ---
the map $T_{s,t}(h):=\delta_{s,t}(\mathbf 1\otimes h)$ --- is bounded
and linear, and satisfies the interval additivity
$T_{r,t}(h_{r,s}\oplus h_{s,t})=T_{r,s}(h_{r,s})+T_{s,t}(h_{s,t})$
inside $L^2(\mathcal F_{r,t})$. Brownian and Poisson divergences both
meet this requirement; in the Brownian case the seed map
$T^W_{s,t}\colon L^2([s,t])\to L^2(\mathcal F_{s,t})$ is the It\^o
isometry, and its range lies in $I_1(\mathcal F_{s,t})$.

The observation of \cref{sec:seed} is that this is no accident. For
\emph{any} family $T_{s,t}$ satisfying interval additivity and
$T_{s,t}(h)\in L^2(\mathcal F_{s,t})$, the elements
$X_{s,t}:=T_{s,t}(h_{s,t})$ form an additive integrable family in
the noise $\mathcal N$, hence belong to the first chaos
$I_1(\mathcal N)$. When $\mathcal N$ is black, $I_1(\mathcal N)$ is
trivial after centering, and $T_{s,t}\equiv 0$. The classical seed
\emph{has no nontrivial instance} over a black noise.

We emphasize what this observation does and does not say. It does not
say that no primitive divergence $\delta_{\mathcal N}$ can exist over
a black noise. It says only that any such $\delta_{\mathcal N}$, if it
exists, cannot be reached by the classical route: the
deterministic-coefficient, interval-additive, first-order seed is
unavailable. Two possibilities remain. \emph{Either} a primitive
divergence over a black noise admits a nonclassical seed ---
non-additive in $(s,t)$, multiplicative, quadratic, or based on
interval functionals of a fundamentally different type --- and an
energy-first It\^o calculus extends beyond classical innovation.
\emph{Or} every candidate construction reduces, after closure, to an
interval-additive seed, in which case the observation forces
triviality and a nontrivial first chaos is a necessary input to the
energy-first program.

\subsection{A second pressure point at the Leibniz rule}\label{sec:intro-leibniz}

The classical Malliavin pair satisfies the Leibniz rule
$D(FG)=F\,DG+G\,DF$, which makes the adjoint $D=\delta^{*}$ a
first-order derivation on its natural test algebra. There is no a
priori reason for this rule to survive when the seed of
$\delta_{\mathcal N}$ ceases to be classical. We record this as a
second, complementary diagnostic in \cref{sec:leibniz}, isolating the
\emph{Leibniz defect}
\[
  \mathfrak L_D(F,G):=D(FG)-F\,DG-G\,DF.
\]
The Wiener case has $\mathfrak L_{D^W}\equiv 0$ on its smooth
cylindrical algebra; the compensated Poisson case has the nonzero
quadratic defect $\mathfrak L_{D^P}(F,G)(z)=(D^P_z F)(D^P_z G)$,
already inside the classical L\'evy world. The diagnostic is that an
energy-first calculus over a nonclassical noise should be expected,
on the Poisson precedent, to have a nonzero Leibniz defect; the
question is whether this defect can be expressed intrinsically from
the noise tensor factorization.

The first pressure point is at the seed; the second is at the
product rule of the adjoint. Together they locate, in
\cref{sec:open}, the central open problem of the paper. The positive
candidate that remains after additive first-chaos seeds vanish is
not absence of structure but a higher-order one: a
\emph{carr\'e-du-champ} / product-defect seed
(\cref{sec:cdc-seed,sec:twopoint}), supplying a positive bilinear
energy kernel in settings where the classical first-chaos seed is
unavailable, and more generally the generator-level interaction
hierarchy of \cref{sec:twopoint-hierarchy}, whose higher members
$\Gamma_{S}$ for $|S|\ge 3$ record the genuinely $|S|$-body
infinitesimal interaction content of the noise above the two-point
layer.

\subsection{What this paper does not do}\label{sec:intro-non-claims}

The framing above is calibrated against several non-claims that we
make explicit here and repeat in the body of the paper.

We do not prove any new black-noise theorem. The proposition of
\cref{sec:proposition} is a restatement of \cite{Tsirelson2004} in
the ladder language of this paper; the examples of black noise are
cited, not reproved. We do not construct $\delta_{\mathcal N}$ on any
nontrivial black noise. The pressure-point observation of
\cref{sec:seed} forbids only the classical seed; it does not exclude
nonclassical seeds, and offers no construction of one. The Leibniz
diagnostic of \cref{sec:leibniz} is exactly that --- a diagnostic.
It does not prove that blackness forces a nonzero Leibniz defect.
We do not claim that IBM is black. The natural noise of IBM is
classical, and the trilogy's obstructions live at the level of
$L^2$ closure and sector summability inside that classical noise,
not at the level of the noise algebra. Black noise is therefore a
\emph{complementary} test of the energy-first stance, addressing
whether the framework can survive without any classical seed at all,
rather than how it survives the loss of $L^2$ closure within a
classical seed.

\subsection{Roadmap}\label{sec:intro-roadmap}

\Cref{sec:definitions} fixes the basic objects: observable It\^o
identity, energy lift, noise in the Tsirelson--Vershik sense,
classical subnoise, black noise, and classical innovation reduction
in the strong sense required for the proposition of
\cref{sec:proposition} to apply.
\Cref{sec:proposition} states and proves the proposition that black
noise forbids classical innovation reduction.
\Cref{sec:calibration} compares the Wiener, IBM, and black-noise
cases at the level of what disappears when the first chaos becomes
trivial. \Cref{sec:seed} states and proves the pressure-point
observation on interval-additive seeds. \Cref{sec:leibniz}
introduces the Leibniz defect and records the complementary
algebraic diagnostic. \Cref{sec:twopoint} computes the two-point
product defect for stochastic flows of kernels and records, in
\cref{sec:twopoint-hierarchy}, the natural all-$|S|$ extension as
the generator-level interaction hierarchy.
\Cref{sec:open} formulates the central open
problem of constructing a primitive divergence $\delta_{\mathcal N}$
over a nonclassical noise, with both possible outcomes stated
symmetrically. \Cref{sec:examples} gathers the known examples of
black noise as citations. \Cref{sec:discussion} returns to the
relation with the trilogy and to the non-claims of this introduction.

%==============================================================
\section{Definitions}\label{sec:definitions}
%==============================================================

This section fixes the objects used throughout the paper: the
observable It\^o identity, the energy / operator-valued lift, the
Tsirelson--Vershik noise and its classical subnoise, black noise,
and the classical innovation reduction in the strong form required
for the proposition of \cref{sec:proposition} to apply. Definitions
are kept to the minimum needed for the lemma of \cref{sec:seed} and
the open problem of \cref{sec:open}; we do not develop Tsirelson's
theory in full generality.

Throughout, $(\Omega,\mathcal F,\Prob)$ is a fixed probability space;
$E$ is a Polish space; all functions and processes are Borel
measurable; martingales are with respect to a specified,
usual-conditions-completed filtration.

\subsection{\texorpdfstring{Observable It\^o identity}{Observable Ito identity}}\label{sec:def-observable}

\begin{definition}[\texorpdfstring{Observable It\^o identity}{Observable Ito identity}]\label{def:observable}
An \emph{observable It\^o identity} is a triple $(X,L,\mathcal T)$ in
which
\begin{itemize}
\item[$(\mathrm i)$] $X=(X_t)_{t\ge 0}$ is a c\`adl\`ag process on a
Polish space $E$, defined on $(\Omega,\mathcal F,\Prob)$;
\item[$(\mathrm{ii})$] $\mathcal T\subset C_b(E)$ is a linear test class;
\item[$(\mathrm{iii})$] $L\colon\mathcal T\to C_b(E)$ is linear; and
\item[$(\mathrm{iv})$] for every $f\in\mathcal T$ the process
\begin{equation}\label{eq:martingale-problem}
  M_t[f]\;:=\;f(X_t)-f(X_0)-\int_0^t Lf(X_s)\,ds,\qquad t\ge 0,
\end{equation}
is a martingale with respect to the natural filtration
$\mathcal F^X_t:=\sigma(X_u:u\le t)$.
\end{itemize}
\end{definition}

\begin{remark}\label{rem:observable-semigroup}
When the martingale problem $(L,\mathcal T)$ is well posed in the
standard sense (see~\cite{EthierKurtz}), the observable It\^o
identity determines the law of $X$, hence its transition semigroup.
It does \emph{not} determine how the randomness in $(M_t[f])$ is
organized as a stochastic process: distinct sources of innovation
can generate identical martingale problems. The present paper is
concerned precisely with the gap between the operator-level identity
\eqref{eq:martingale-problem} and any further representation of
$(M_t[f])$ in terms of an underlying integrator.
\end{remark}

\subsection{Energy lift}\label{sec:def-energy}

\begin{definition}[Energy lift]\label{def:energy-lift}
An \emph{energy lift} of an observable It\^o identity
$(X,L,\mathcal T)$ is a quintuple
\[
  \bigl(\mathcal H,\;\mathcal U,\;\delta,\;D,\;(u^{\,\cdot}_f)_{f\in\mathcal T}\bigr)
\]
in which:
\begin{itemize}
\item[$(\mathrm i)$] $\mathcal H$ is a complex Hilbert space, or a
rigged Hilbert triple $\mathcal H_+\subset\mathcal H\subset\mathcal H_-$,
with $L^2(\Omega,\mathcal F^X,\Prob)\subset\mathcal H$ as a dense
subspace;
\item[$(\mathrm{ii})$] $\mathcal U$ is a Hilbert space of test directions;
\item[$(\mathrm{iii})$] $\delta\colon\Dom(\delta)\subset\mathcal U\to L^2(\Omega)$
is a closed, densely-defined operator (the \emph{divergence}), with
closed adjoint $D=\delta^{*}\colon\Dom(D)\subset L^2(\Omega)\to\mathcal U$
(the \emph{adjoint derivative}), satisfying the integration-by-parts
identity
\begin{equation}\label{eq:IBP}
  \E\!\bigl[F\cdot\delta(u)\bigr]
  \;=\;\E\!\bigl[\langle DF,u\rangle_{\mathcal U}\bigr],
  \qquad F\in\Dom(D),\ u\in\Dom(\delta);
\end{equation}
\item[$(\mathrm{iv})$] for every $f\in\mathcal T$ there is a
predictable direction
$u^{\,\cdot}_f\colon[0,\infty)\to\mathcal U$ such that
\begin{equation}\label{eq:M-as-delta}
  M_t[f]\;=\;\delta\!\bigl(u_f^{[0,t]}\bigr)\quad\text{in }L^2(\Omega),
  \qquad t\ge 0,
\end{equation}
where $u_f^{[0,t]}$ is the restriction of $u_f^{\,\cdot}$ to $[0,t]$.
\end{itemize}
\end{definition}

\begin{remark}\label{rem:energy-lift-no-classical}
\Cref{def:energy-lift} makes no commitment as to the structure of
$\delta$. In particular, $\delta$ is not assumed to be the
It\^o--Wiener integral against a Brownian motion, nor the
compensated integral against a Poisson random measure, nor any other
classical L\'evy-type integrator. The Wiener and Poisson divergences
are particular realizations of $\delta$; the energy-first stance of
this paper treats $\delta$ as the primitive object and asks under
what conditions on the underlying noise such a $\delta$ can exist
at all. The terminology \emph{adjoint derivative} for $D=\delta^{*}$
is deliberate: the Leibniz rule $D(FG)=F\,DG+G\,DF$, and hence the
property of being a derivation on a test algebra, is an additional
optional structure rather than part of the basic definition. See
\cref{sec:leibniz}.
\end{remark}

\subsection{Noise}\label{sec:def-noise}

\begin{definition}[Noise]\label{def:noise}
A \emph{noise} on $(\Omega,\mathcal F,\Prob)$ is a system
$\mathcal N=\bigl((\mathcal F_{s,t})_{s\le t},(T_h)_{h\in\R}\bigr)$
of sub-$\sigma$-fields and a measurable measure-preserving group
action, satisfying:
\begin{itemize}
\item[$(\mathrm i)$] \emph{Independent factorization.} For every
$r\le s\le t$,
\begin{equation}\label{eq:noise-factor}
  \mathcal F_{r,t}\;=\;\mathcal F_{r,s}\vee\mathcal F_{s,t},
  \qquad \mathcal F_{r,s}\perp\!\!\!\perp\mathcal F_{s,t}.
\end{equation}
\item[$(\mathrm{ii})$] \emph{Stationarity.} Each
$T_h\colon\Omega\to\Omega$ is measure preserving,
$T_{h_1+h_2}=T_{h_1}\circ T_{h_2}$, and
$T_h(\mathcal F_{s,t})=\mathcal F_{s+h,\,t+h}$.
\item[$(\mathrm{iii})$] \emph{Continuity.} The standard continuity
conditions on $(\mathcal F_{s,t})$ in $(s,t)$ hold (right-continuity
at right endpoints; left-continuity at left endpoints).
\end{itemize}
We refer to~\cite{Tsirelson2004}, \S 2, for the full
Tsirelson--Vershik axioms; the form above is sufficient for the
present paper.
\end{definition}

\begin{remark}[\texorpdfstring{$L^2$-tensor form}{L^2-tensor form}]\label{rem:noise-l2-tensor}
The independent factorization \eqref{eq:noise-factor} is equivalent,
after passing to $L^2$, to the unitary tensor identification
\begin{equation}\label{eq:L2-tensor}
  L^2(\mathcal F_{r,t})\;\cong\;L^2(\mathcal F_{r,s})\otimes L^2(\mathcal F_{s,t}),
  \qquad r\le s\le t,
\end{equation}
under which the constants $\mathbf 1$ correspond to the canonical
embeddings. In \cref{sec:seed} we write
$\mathcal F_{r,t}=\mathcal F_{r,s}\otimes\mathcal F_{s,t}$ as shorthand
for the $L^2$-level identification \eqref{eq:L2-tensor}; we never use
$\otimes$ as an algebraic tensor product of $\sigma$-fields.
\end{remark}

\subsection{First chaos and classical subnoise}\label{sec:def-classical}

\begin{definition}[Additive integrable family]\label{def:additive-family}
Let $\mathcal N=(\mathcal F_{s,t})$ be a noise. An \emph{additive
integrable family} on $\mathcal N$ is a collection
$X=(X_{s,t})_{s\le t}$ such that:
\begin{itemize}
\item[$(\mathrm i)$] $X_{s,t}\in L^2(\mathcal F_{s,t})$ for every $s\le t$;
\item[$(\mathrm{ii})$] for every $r\le s\le t$,
\begin{equation}\label{eq:additivity}
  X_{r,t}\;=\;X_{r,s}+X_{s,t}\quad\text{in }L^2(\mathcal F_{r,t}),
\end{equation}
where the right-hand side is read via the canonical embeddings
$L^2(\mathcal F_{r,s})\hookrightarrow L^2(\mathcal F_{r,t})$ and
$L^2(\mathcal F_{s,t})\hookrightarrow L^2(\mathcal F_{r,t})$ of
\cref{rem:noise-l2-tensor}.
\end{itemize}
The family $X$ is \emph{centered} if $\E[X_{s,t}]=0$ for every $s\le t$.
\end{definition}

\begin{definition}[First chaos; classical subnoise]\label{def:classical-subnoise}
The \emph{first chaos} $I_1(\mathcal N)$ of $\mathcal N$ is the
closed linear subspace of $\bigoplus_{s\le t}L^2(\mathcal F_{s,t})$
spanned by all additive integrable families on $\mathcal N$ in the
sense of \cref{def:additive-family}. The \emph{classical subnoise}
$\mathcal N^{\mathrm{cl}}=(\mathcal F^{\mathrm{cl}}_{s,t})$ is the
subnoise generated by $I_1(\mathcal N)$:
\[
  \mathcal F^{\mathrm{cl}}_{s,t}
  :=\sigma\bigl(\,X_{s,t}\;:\;X\text{ an additive integrable family on }\mathcal N\,\bigr),
  \qquad s\le t.
\]
\end{definition}

We refer to~\cite{Tsirelson2004}, \S 3, for the development of $I_1$
and $\mathcal N^{\mathrm{cl}}$, and for their basic structural
properties (Hilbertian closure, behavior under refinement,
equivalence with the \emph{stable} or \emph{linearizable} subnoise).

\subsection{Black noise}\label{sec:def-black}

\begin{definition}[Black noise]\label{def:black}
A noise $\mathcal N$ is \emph{black} if its classical subnoise
$\mathcal N^{\mathrm{cl}}$ is trivial,
\[
  \mathcal F^{\mathrm{cl}}_{s,t}\;=\;\{\emptyset,\Omega\}\bmod\Prob
  \qquad\text{for every }s<t,
\]
while $\mathcal N$ itself is nontrivial in the sense that
$\mathcal F_{s,t}\neq\{\emptyset,\Omega\}\bmod\Prob$ for some $s<t$.
\end{definition}

\begin{remark}[Equivalent form]\label{rem:black-no-centered-additive}
By \cref{def:additive-family} and \cref{def:classical-subnoise},
$\mathcal N$ is black if and only if every centered additive
integrable family $(X_{s,t})$ on $\mathcal N$ satisfies $X_{s,t}=0$
almost surely for every $s\le t$. In this equivalent form, blackness
asserts the absence of any nonconstant first-order, additive
$L^2$-content in $\mathcal N$. Known examples of black noise include
the noise of Tsirelson's sticky and coalescing
flows~\cite{LeJanRaimond2004}, the noise of the Brownian
web~\cite{Tsirelson2004,SSS2017}, and the noise of two-dimensional
critical percolation~\cite{SchrammSmirnov2011}; we cite these in
\cref{sec:examples} and do not reprove them here.
\end{remark}

\subsection{Natural noise; classical innovation reduction}\label{sec:def-cir}

\begin{definition}[Natural noise of a process]\label{def:natural-noise}
Let $(X,L,\mathcal T)$ be an observable It\^o identity. We assume
that $X$ carries a \emph{natural noise}
$\mathcal N_X=(\mathcal F^X_{s,t})_{s\le t}$ in the sense of
\cref{def:noise}, with $\mathcal F^X_{s,t}\subset\mathcal F$ for every
$s\le t$ and
$\bigvee_{u\le t}\mathcal F^X_{0,u}=\mathcal F^X_t$. For processes
with stationary, independent-like increments --- Brownian, L\'evy,
Poisson; iterated and stochastic-flow examples treated in
\cref{sec:calibration,sec:examples} --- $\mathcal N_X$ may be taken
as the natural Tsirelson--Vershik noise generated by $X$. For
general $X$ the existence of $\mathcal N_X$ is an assumption on the
model, not a universal construction. We write
$\mathcal N_X^{\mathrm{cl}}$ for the classical subnoise of
$\mathcal N_X$.
\end{definition}

\begin{definition}[Classical innovation reduction]\label{def:cir}
A \emph{classical innovation reduction} of an observable It\^o
identity $(X,L,\mathcal T)$ is a stochastic process
$Z=(Z_t)_{t\ge 0}$ on $(\Omega,\mathcal F,\Prob)$ such that:
\begin{itemize}
\item[$(\mathrm i)$] $Z$ is a L\'evy process (Brownian, Poisson,
compound Poisson, or finite-dimensional L\'evy), and is adapted to
the natural filtration $\mathcal F^X_t$;
\item[$(\mathrm{ii})$] the noise $\mathcal N_Z$ generated by $Z$ is a
subnoise of the classical subnoise $\mathcal N_X^{\mathrm{cl}}$ of
$X$, and is nontrivial in the sense that
$\mathcal F^Z_{s,t}\neq\{\emptyset,\Omega\}\bmod\Prob$ for some $s<t$;
\item[$(\mathrm{iii})$] for every $f\in\mathcal T$ there exists an
$\mathcal F^X$-predictable integrand $H[f]$ such that
\begin{equation}\label{eq:cir-representation}
  M_t[f]\;=\;\int_0^t H_s[f]\,dZ_s\quad\text{in }L^2(\Omega),
  \qquad t\ge 0,
\end{equation}
and the predictable representation $H[f]$ determines $M[f]$ uniquely.
\end{itemize}
\end{definition}

\begin{remark}[Adaptedness and the enlargement caveat]\label{rem:cir-no-enlargement}
Condition $(\mathrm i)$ is what makes \cref{def:cir} bite. Adjoining
a Brownian motion $W'$ to $(\Omega,\mathcal F,\Prob)$ on an enlarged
probability space, independent of all existing randomness, and
writing $M_t[f]=\int_0^t H_s[f]\,dW'_s$ for some predictable $H[f]$
on the enlargement does \emph{not} produce a classical innovation
reduction in the present sense, because $W'$ is not adapted to
$\mathcal F^X$. Condition $(\mathrm{ii})$ further requires that the
L\'evy driver contribute to the classical part of the noise of $X$.
In combination, $(\mathrm i)$--$(\mathrm{ii})$ formalize the
requirement that the L\'evy representation
\eqref{eq:cir-representation} take place \emph{inside} the noise
$\mathcal N_X$ rather than on an extraneous enlargement.
\end{remark}

\begin{remark}[Strength relative to the energy lift]\label{rem:cir-vs-energy-lift}
\Cref{def:cir} is strictly stronger than \cref{def:energy-lift}.
Every classical innovation reduction yields an energy lift, with
$\delta=\delta^Z$ the It\^o--Wiener (or compensated Poisson)
divergence against $Z$ and $D=D^Z$ its Malliavin adjoint. The
converse fails: an energy lift may exist without any $Z$ satisfying
\cref{def:cir}. The proposition of \cref{sec:proposition} records the
precise mechanism by which blackness of $\mathcal N_X$ forbids the
converse.
\end{remark}

%==============================================================
\section{Classical innovation reduction and black noise}\label{sec:proposition}
%==============================================================

We record the benchmark proposition showing that blackness of
$\mathcal N_X$ forbids the classical innovation reduction of
\cref{def:cir}. The statement is short and conservative:
a restatement of Tsirelson's blackness in the form required by the
operator-derivative framework, not a new black-noise theorem.

\begin{proposition}\label{prop:no-cir}
Let $(X,L,\mathcal T)$ be an observable It\^o identity with natural
noise $\mathcal N_X$ in the sense of \cref{def:natural-noise}.
Suppose $\mathcal N_X$ is black in the sense of \cref{def:black}.
Then $(X,L,\mathcal T)$ admits no classical innovation reduction in
the sense of \cref{def:cir}: no process $Z$ satisfying conditions
$(\mathrm i)$--$(\mathrm{iii})$ of \cref{def:cir} can represent the
martingale family $(M_t[f])$ as in
\eqref{eq:cir-representation}.
\end{proposition}

\begin{proof}
By \cref{def:cir}$(\mathrm{ii})$, a classical innovation reduction
requires the noise $\mathcal N_Z$ generated by the candidate driver
$Z$ to be a \emph{nontrivial} subnoise of the classical subnoise
$\mathcal N_X^{\mathrm{cl}}$. Blackness of $\mathcal N_X$
(\cref{def:black}) means $\mathcal N_X^{\mathrm{cl}}$ is trivial, so
every subnoise of $\mathcal N_X^{\mathrm{cl}}$ is trivial. This
contradicts the required nontriviality of $\mathcal N_Z$, and no
classical innovation reduction exists.

Equivalently, any L\'evy driver whose interval noise genuinely lives
inside the natural noise of $X$ would generate nontrivial classical
additive $L^2$-content on $\mathcal N_X$, which blackness forbids.
\end{proof}

\begin{remark}\label{rem:prop-tsirelson-restatement}
\Cref{prop:no-cir} is essentially a restatement of the definition of
black noise in the ladder language of this paper. The argument is
direct from the definitions: condition $(\mathrm{ii})$ of
\cref{def:cir} demands that the noise of the candidate L\'evy driver
be a nontrivial subnoise of $\mathcal N_X^{\mathrm{cl}}$, and
blackness of $\mathcal N_X$ collapses that classical subnoise to a
trivial one. We do not advertise this as a new black-noise theorem.
Its role here is to make precise, in the framework of
\cref{def:energy-lift,def:cir}, what blackness obstructs and what it
does not.
\end{remark}

\begin{remark}[What is not asserted]\label{rem:prop-non-claim}
\Cref{prop:no-cir} asserts only that the strong classical reduction
of \cref{def:cir} is unavailable when $\mathcal N_X$ is black. It
does not assert that no energy lift of $(X,L,\mathcal T)$ in the
sense of \cref{def:energy-lift} exists, and it does not assert that
no primitive divergence $\delta_{\mathcal N}$ over $\mathcal N_X$
exists. Those are exactly the questions to which
\cref{sec:seed,sec:leibniz,sec:open} return.
\end{remark}

%==============================================================
\section{Calibration: Wiener, IBM, black noise}\label{sec:calibration}
%==============================================================

Before turning to the technical pressure points of
\cref{sec:seed,sec:leibniz}, it is useful to set the trilogy and the
black-noise setting side by side at the level of what the first
chaos provides. We summarize this in three short paragraphs and one
metaphor; nothing in this section is proved here.

\subsection{Wiener noise}\label{sec:cal-wiener}

Let $W$ be a standard Brownian motion. Its noise $\mathcal N_W$ is
classical and the first chaos
$I_1(\mathcal F^W_{s,t})\cong L^2([s,t])$ is canonically Hilbertian:
each $h\in L^2([s,t])$ corresponds to $\int_s^t h(u)\,dW_u\in
I_1(\mathcal F^W_{s,t})$. The increment seed $W_t-W_s$ is the
prototypical centered additive integrable family. From this seed the
It\^o divergence $\delta^W$ and its Malliavin adjoint $D^W$ are
built by closure and chaos extension; the Wiener case is the source
of the classical pattern that the energy-first program generalizes.

\subsection{IBM and the trilogy}\label{sec:cal-ibm}

Iterated Brownian motion $\IBM_t:=X_{|Y_t|}$, with $X$ a two-sided
Brownian scenery and $Y$ an independent Brownian clock, lies on a
classical noise $\mathcal N_{(X,Y)}=\mathcal N_X\otimes\mathcal N_Y$
in the sense of~\eqref{eq:L2-tensor}. Both factors carry rich first
chaoses. The trilogy's obstructions arise \emph{inside} this
classical setting and are not filtration-level:
\begin{itemize}
\item[--] \cite{EnergyFirst} mollifies the scenery,
$X^\varepsilon=\rho_\varepsilon*X$, and proves an energy-first It\^o
formula at the mollified level; the natural $L^2$ closure of the
correction term fails on a first-$X$-chaos sector, and only an
$H^{-1}$ dual closure survives;
\item[--] \cite{FCDC} identifies the dual closure of that
first-chaos sector explicitly;
\item[--] \cite{BidegreeAtlas} organizes the higher bidegree sectors
$(m,n)$ along the chaos decompositions of $\mathcal N_X$ and
$\mathcal N_Y$ into an atlas with their own summability and closure
conditions.
\end{itemize}
None of these obstructions is at the level of the first chaos or the
noise algebra: the seed of the classical construction is intact,
and the failure occurs in the closure (\cite{EnergyFirst,FCDC}) and
the cross-sector summability (\cite{BidegreeAtlas}).

\subsection{Black noise}\label{sec:cal-black}

For a black noise $\mathcal N$, the first chaos collapses:
$I_1(\mathcal N)$ is trivial modulo constants, by
\cref{def:black,rem:black-no-centered-additive}. The Brownian
template breaks at the first step. There is no $\mathbf 1_{[s,t]}$
to integrate against a classical generator, no isometry
$L^2([s,t])\to I_1(\mathcal F_{s,t})$ to extend by closure, and no
chaos decomposition $\bigoplus_n I_n$ built from tensor powers of
$I_1$. Any candidate energy-first construction must therefore
originate elsewhere, at the algebra of the noise rather than at a
first-chaos generator.

\subsection{Horizontal and vertical stresses}\label{sec:cal-metaphor}

The trilogy and the black-noise setting stress the energy-first
program along orthogonal axes. The trilogy operates \emph{within}
the classical first-chaos input and asks how far the resulting
calculus can be carried when its natural $L^2$ closures fail and
when the bidegree sectors must be summed nonuniformly; this is a
\emph{horizontal} stress test. Black noise operates \emph{below} the
first chaos and asks whether the calculus can be carried out when
no classical first-chaos input is available at all; this is a
\emph{vertical} stress test. The two settings are complementary: a
positive resolution of either does not entail the other, and the
non-trivial obstructions in the two cases are of different types.
The trilogy is not concerned with blackness, and the present paper
is not concerned with $L^2$-closure failures inside classical noise.

%==============================================================
\section{The first structural pressure point}\label{sec:seed}
%==============================================================

The classical It\^o construction of a divergence operator begins not
with a closed unbounded operator on Fock space but with a much
smaller object: a bounded linear map from deterministic test
functions on an interval into the first chaos of the noise on that
interval. Closure, chaos extension, and the full Malliavin pair are
obtained from this seed by standard limiting arguments. In this
section we record the precise structural form of the seed in the
classical case and observe that, over a black noise, no nontrivial
seed of this form exists. The observation is short and conservative;
it locates the first place where the classical template fails,
without claiming that the failure is fatal to the energy-first
program.

\subsection{Setup}\label{sec:seed-setup}

Let $\mathcal N=(\mathcal F_{s,t})_{s\le t}$ be a noise in the
Tsirelson--Vershik sense (\cref{def:noise}), with the tensor
factorization
\begin{equation}\label{eq:tensor-factorization}
  \mathcal F_{r,t}=\mathcal F_{r,s}\otimes\mathcal F_{s,t},
  \qquad r\le s\le t,
\end{equation}
read at the $L^2$ level via~\eqref{eq:L2-tensor}, and the associated
unitary identification
\[
  L^2(\mathcal F_{r,t})\;\cong\;L^2(\mathcal F_{r,s})\otimes L^2(\mathcal F_{s,t}).
\]
We write $\iota^{r,t}_{s,t}\colon L^2(\mathcal F_{s,t})\hookrightarrow L^2(\mathcal F_{r,t})$
for the canonical embedding $X\mapsto\mathbf 1_{r,s}\otimes X$, and
$\iota^{r,t}_{r,s}$ for the symmetric embedding
$X\mapsto X\otimes\mathbf 1_{s,t}$. We suppress these embeddings
when no ambiguity arises.

Let $(H_{s,t})_{s\le t}$ be a family of separable Hilbert spaces of
\emph{deterministic test directions}, with an additive structure
\begin{equation}\label{eq:H-additive}
  H_{r,t}\;=\;H_{r,s}\oplus H_{s,t},\qquad r\le s\le t,
\end{equation}
and isometric inclusions $H_{s,t}\hookrightarrow H_{r,t}$ compatible
with refinement of intervals. The canonical example is
$H_{s,t}=L^2([s,t])$.

\begin{definition}[Interval-additive deterministic seed]\label{def:seed}
An \emph{interval-additive deterministic seed} on $\mathcal N$,
valued in $(H_{s,t})$, is a family
\[
  T_{s,t}\colon H_{s,t}\longrightarrow L^2(\mathcal F_{s,t}),
  \qquad s\le t,
\]
of bounded linear maps satisfying, for every $r\le s\le t$ and every
$h_{r,s}\in H_{r,s}$, $h_{s,t}\in H_{s,t}$,
\begin{equation}\label{eq:interval-additive}
  T_{r,t}(h_{r,s}\oplus h_{s,t})
  \;=\;
  \iota^{r,t}_{r,s}\!\bigl(T_{r,s}(h_{r,s})\bigr)
  +\iota^{r,t}_{s,t}\!\bigl(T_{s,t}(h_{s,t})\bigr)
  \quad\text{in }L^2(\mathcal F_{r,t}).
\end{equation}
\end{definition}

The definition isolates the structural content of the classical seed.
If $\delta_{\mathcal N}$ is assembled from a family of operators
$\delta_{s,t}\colon\Dom(\delta_{s,t})\subset L^2(\mathcal F_{s,t};H_{s,t})\to L^2(\mathcal F_{s,t})$
as in the open problem of \cref{sec:open}, and one restricts to
deterministic test coefficients by setting
\begin{equation}\label{eq:seed-from-delta}
  T_{s,t}(h)\;:=\;\delta_{s,t}(\mathbf 1\otimes h),
  \qquad h\in H_{s,t},
\end{equation}
then $T_{s,t}$ inherits boundedness from a basic continuity
assumption on $\delta_{s,t}$ and inherits
\eqref{eq:interval-additive} from any reasonable splitting of
$\delta_{r,t}$ along \eqref{eq:tensor-factorization}. We call
\eqref{eq:seed-from-delta} the \emph{deterministic-coefficient seed}
of $\delta_{\mathcal N}$.

\subsection{Brownian and Poisson benchmarks}\label{sec:seed-benchmarks}

Let $W$ be a standard Brownian motion on $[0,\infty)$ with natural
noise $\mathcal N_W=(\mathcal F^W_{s,t})$, and take $H_{s,t}=L^2([s,t])$
with the additive structure~\eqref{eq:H-additive} provided by
$L^2([r,t])=L^2([r,s])\oplus L^2([s,t])$. The Wiener seed
\begin{equation}\label{eq:wiener-seed}
  T^W_{s,t}(h)\;:=\;\int_s^t h(u)\,dW_u
\end{equation}
is bounded linear, $\|T^W_{s,t}(h)\|_{L^2}=\|h\|_{L^2([s,t])}$, and
satisfies \eqref{eq:interval-additive} by the standard splitting of
the It\^o integral. The range of $T^W_{s,t}$ is exactly the first
chaos $I_1(\mathcal F^W_{s,t})$, and $T^W$ is precisely the seed from
which the Malliavin divergence $\delta^W$ is reconstructed by chaos
extension.

For a Poisson process $N$ of intensity $\lambda$ with compensated
version $\widetilde N_t=N_t-\lambda t$, the seed
$T^{\widetilde N}_{s,t}(h)=\int_s^t h(u)\,d\widetilde N_u$ is bounded
linear, satisfies the analogous splitting, and ranges into the
compensated first chaos of $\mathcal N_{\widetilde N}$. In both
cases the seed is the input from which the full divergence operator
is built by closure and chaos extension; in both cases it satisfies
\cref{def:seed} and lands in $I_1$.

\subsection{First-chaos confinement}\label{sec:seed-lemma}

We now record the structural observation. Let
$H_\infty:=\bigcup_{s\le t}H_{s,t}$ be the directed union under
refinement, and call $h\in H_\infty$ \emph{compatible} if it has, for
every interval $(s,t)$, a component $h_{s,t}\in H_{s,t}$ such that
$h_{r,t}=h_{r,s}\oplus h_{s,t}$ under~\eqref{eq:H-additive}. In the
Brownian benchmark this is exactly $h\in L^2([0,T])$ with
$h_{s,t}=h\,\mathbf 1_{[s,t]}$.

\begin{lemma}[First-chaos confinement]\label{lem:seed}
Let $\mathcal N=(\mathcal F_{s,t})$ be a noise and let $(T_{s,t})$ be
an interval-additive deterministic seed in the sense of
\cref{def:seed}. For every compatible $h\in H_\infty$, define
\begin{equation}\label{eq:X-family}
  X_{s,t}\;:=\;T_{s,t}(h_{s,t})-\E\!\bigl[T_{s,t}(h_{s,t})\bigr],
  \qquad s\le t.
\end{equation}
Then:
\begin{enumerate}
\item[$(\mathrm i)$] $X_{s,t}\in L^2(\mathcal F_{s,t})$ is centered
for every $s\le t$;
\item[$(\mathrm{ii})$] $(X_{s,t})$ is additive:
\[
  X_{r,t}\;=\;X_{r,s}+X_{s,t}
  \quad\text{in }L^2(\mathcal F_{r,t}),\quad r\le s\le t;
\]
\item[$(\mathrm{iii})$] $(X_{s,t})\in I_1(\mathcal N)$.
\end{enumerate}
If, in addition, $\mathcal N$ is black in the sense of
\cref{def:black}, then $X_{s,t}=0$ almost surely for every $s\le t$,
and the seed reduces to its deterministic mean:
\begin{equation}\label{eq:seed-deterministic}
  T_{s,t}(h_{s,t})\;=\;\E\!\bigl[T_{s,t}(h_{s,t})\bigr]
  \qquad\text{for every compatible }h.
\end{equation}
\end{lemma}

\begin{proof}
Claim $(\mathrm i)$ is immediate: $T_{s,t}$ maps into
$L^2(\mathcal F_{s,t})$ by \cref{def:seed}, and centering is by
construction.

For $(\mathrm{ii})$, apply \eqref{eq:interval-additive} with
$h_{r,t}=h_{r,s}\oplus h_{s,t}$ and take expectations:
\begin{align*}
  X_{r,t}
  &=T_{r,t}(h_{r,t})-\E[T_{r,t}(h_{r,t})]\\
  &=T_{r,s}(h_{r,s})+T_{s,t}(h_{s,t})
   -\E[T_{r,s}(h_{r,s})]-\E[T_{s,t}(h_{s,t})]\\
  &=X_{r,s}+X_{s,t},
\end{align*}
with the embeddings $\iota$ suppressed.

For $(\mathrm{iii})$, $(X_{s,t})$ is by $(\mathrm i)$--$(\mathrm{ii})$
a centered, $\mathcal F_{s,t}$-measurable, square-integrable, additive
family on $\mathcal N$. By \cref{def:additive-family} and
\cref{def:classical-subnoise}, $I_1(\mathcal N)$ is the closed
linear span in $L^2$ of all such families, so
$(X_{s,t})\in I_1(\mathcal N)$.

The final assertion is immediate. Blackness of $\mathcal N$ means
$\mathcal N^{\mathrm{cl}}$ is trivial, equivalently
$I_1(\mathcal N)=\{0\}$ modulo constants
(\cref{rem:black-no-centered-additive}). The centering in
$(\mathrm i)$ removes the constant part, so $X_{s,t}=0$ in $L^2$,
hence almost surely. This is \eqref{eq:seed-deterministic}.
\end{proof}

\subsection{Interpretation}\label{sec:seed-interpretation}

\Cref{lem:seed} says, with no hypothesis on $\mathcal N$ beyond the
noise axioms, that the centered part of any interval-additive
deterministic seed lives in the first chaos. The Wiener and Poisson
seeds saturate this confinement: their ranges are precisely the
additive integrable families generating $I_1$. In a black noise,
where $I_1$ is trivial, the same confinement forces the centered seed
to vanish entirely.

It is essential to be clear about the scope of this observation.

\medskip
\noindent\emph{What the lemma does say.}
The classical route to building $\delta_{\mathcal N}$ --- starting
from a bounded, linear, deterministic-coefficient family of interval
operators $T_{s,t}\colon H_{s,t}\to L^2(\mathcal F_{s,t})$ that is
additive under~\eqref{eq:H-additive} --- is unavailable over any
black noise. The seed, if assembled in this manner, is forced to be
deterministic and therefore carries no stochastic content. The
Wiener template $T^W_{s,t}(h)=\int_s^t h\,dW$ has no nonzero analogue
over a black noise, because the very $L^2([s,t])$-isometric image
into $I_1$ that underwrites~\eqref{eq:wiener-seed} does not exist
there.

\medskip
\noindent\emph{What the lemma does not say.}
\Cref{lem:seed} does not exclude the existence of a primitive
divergence $\delta_{\mathcal N}$ over a black noise. It constrains
only one form of seed --- the bounded linear, interval-additive,
deterministic-coefficient form. Several escape routes remain open
and are not addressed by the lemma:
\begin{itemize}
\item[$(\mathrm a)$] \emph{Non-additive interval structure.} The
hypothesis~\eqref{eq:H-additive} on test-direction spaces is what
feeds additivity into the seed. A noise-intrinsic test class with a
multiplicative or otherwise non-additive composition law on
intervals is not constrained by \cref{lem:seed}.
\item[$(\mathrm b)$] \emph{Nonlinear or higher-order seeds.} A seed
that is quadratic in $h$, or of area- or bracket-type in the sense
of rough paths, produces on each interval not an additive family
$(X_{s,t})$ but a higher-order interval object; such objects need
not lie in $I_1$.
\item[$(\mathrm c)$] \emph{Stochastic coefficients.} \Cref{lem:seed}
constrains only the deterministic-coefficient seed
$T_{s,t}(h)=\delta_{s,t}(\mathbf 1\otimes h)$. A primitive divergence
defined directly on stochastic integrands, without first specializing
to deterministic test coefficients, is not constrained.
\item[$(\mathrm d)$] \emph{Domain choice.} The dense domain in
$L^2(\mathcal F_{s,t})$ on which $\delta_{s,t}$ acts is part of the
data; the lemma constrains only the deterministic-coefficient slice
of that domain.
\end{itemize}

\begin{remark}\label{rem:seed-interpretation}
\Cref{lem:seed} is \emph{not} an impossibility theorem for a
primitive divergence $\delta_{\mathcal N}$, and not an impossibility
theorem for a Malliavin-type calculus over a black noise. It is a
pressure-point observation: it identifies precisely which classical
move is unavailable, and locates the first nontrivial design
decision that any candidate construction must make. The fork
produced by the lemma --- either $\delta_{\mathcal N}$ admits a
nonclassical seed in one of the senses $(\mathrm a)$--$(\mathrm d)$,
or no nontrivial $\delta_{\mathcal N}$ exists over a black noise ---
is the question to which \cref{sec:open} returns, and which the
present paper leaves open.
\end{remark}

%==============================================================
\section{Leibniz defect and nonclassical seeds}\label{sec:leibniz}
%==============================================================

The seed lemma of \cref{sec:seed} locates the first structural
pressure point at the level of the deterministic-coefficient,
interval-additive seed of $\delta_{\mathcal N}$. A second,
complementary pressure point lies one level higher in the calculus,
in the algebraic structure of the adjoint derivative
$D=\delta^{*}$. The classical Malliavin pair satisfies the Leibniz
rule, which makes $D$ a first-order derivation on its natural test
algebra. There is no a priori reason for this rule to survive when
the seed of $\delta_{\mathcal N}$ ceases to be classical; in this
section we record the natural diagnostic and the open question that
goes with it.

\subsection{Test algebra and Leibniz defect}\label{sec:leibniz-defect}

Throughout this section we assume that, in addition to the data of
\cref{def:energy-lift}, there are given:
\begin{itemize}
\item[$(\mathrm a)$] a test algebra
$\mathcal A\subset L^\infty(\Omega)\cap\Dom(D)$, closed under
pointwise products: $FG\in\mathcal A$ for every $F,G\in\mathcal A$;
\item[$(\mathrm b)$] an $L^\infty(\Omega)$-module structure on
$\mathcal U$ extending its Hilbert structure, so that for
$F\in L^\infty(\Omega)$ and $u\in\mathcal U$ the product
$Fu\in\mathcal U$ is well defined and the resulting bilinear map
$L^\infty(\Omega)\times\mathcal U\to\mathcal U$ is bounded.
\end{itemize}
Under $(\mathrm a)$--$(\mathrm b)$, the elements $F\,DG$ and $G\,DF$
lie in $\mathcal U$ for $F,G\in\mathcal A$, and the following object
is well defined.

\begin{definition}[Leibniz defect]\label{def:leibniz-defect}
The \emph{Leibniz defect} of $D$ on $\mathcal A$ is the bilinear map
$\mathfrak L_D\colon\mathcal A\times\mathcal A\to\mathcal U$ given by
\begin{equation}\label{eq:leibniz-defect}
  \mathfrak L_D(F,G)\;:=\;D(FG)-F\,DG-G\,DF,
  \qquad F,G\in\mathcal A.
\end{equation}
We say $D$ is \emph{first-order on $\mathcal A$} if
$\mathfrak L_D\equiv 0$, and that $(\mathcal A,D)$ is a
\emph{derivation pair} in this case.
\end{definition}

\begin{remark}\label{rem:why-adjoint-derivative}
The condition $\mathfrak L_D\equiv 0$ is, by
\cref{def:leibniz-defect}, the condition that $D$ be a derivation
on $\mathcal A$. The general \cref{def:energy-lift} of an energy
lift does not impose this condition: the integration-by-parts
identity~\eqref{eq:IBP} fixes only the closed adjoint pair
$(\delta,D)$ on a dense common domain. The terminology
\emph{adjoint derivative} for $D=\delta^{*}$ is intended to leave
the Leibniz rule, and hence the derivation property, as an
additional optional structure rather than as part of the basic
definition. The classical Malliavin examples below show that the
derivation property is the natural, but not the only, possibility.
\end{remark}

\subsection{Wiener and Poisson benchmarks}\label{sec:leibniz-benchmarks}

Two classical examples show that $\mathfrak L_D\equiv 0$ and
$\mathfrak L_D\not\equiv 0$ both occur naturally inside the
classical L\'evy world.

\emph{Wiener / Malliavin.} Let $W$ be a Brownian motion, $D^W$ the
Malliavin derivative, $\mathcal U=L^2(\R_+)$, and $\mathcal A^W$ the
algebra of smooth cylindrical functionals of $W$. Then
\begin{equation}\label{eq:leibniz-wiener}
  D^W(FG)\;=\;F\,D^W G+G\,D^W F,\qquad F,G\in\mathcal A^W,
\end{equation}
so $\mathfrak L_{D^W}\equiv 0$ on $\mathcal A^W$
(\cite{Nualart}, Prop.~1.2.3). The Malliavin pair $(\delta^W,D^W)$
is a derivation pair.

\emph{Poisson / finite difference.} Let $N$ be a Poisson random
measure on a measurable space $Z$ with intensity $\mu$, let
$\widetilde N=N-\mu$ be its compensated version, and let $D^P$ be the
difference operator
\[
  (D^P_z F)(\omega)\;:=\;F(\omega+\delta_z)-F(\omega),\qquad z\in Z,
\]
extended to a closed operator
$L^2(\Omega)\to L^2(\Omega\times Z,\mu)$ (\cite{Nualart}, \S 6.2;
\cite{LastPenrose}). A direct computation on the natural test algebra
$\mathcal A^P$ gives
\begin{equation}\label{eq:leibniz-poisson}
  D^P_z(FG)
  \;=\;F\,D^P_z G+G\,D^P_z F+(D^P_z F)(D^P_z G),
\end{equation}
so the Leibniz defect of $D^P$ is the nonzero quadratic correction
\begin{equation}\label{eq:leibniz-defect-poisson}
  \mathfrak L_{D^P}(F,G)(z)\;=\;(D^P_z F)(D^P_z G).
\end{equation}
A nonzero defect is therefore not pathological: it is the algebraic
signature of a calculus that is not first-order, exactly as expected
for jump-type integrators.

\subsection{The diagnostic over black noise}\label{sec:leibniz-black}

The natural diagnostic over a black noise combines the seed
confinement of \cref{sec:seed} with the Leibniz framework above. The
proposition below is the simplest version: it says only what
follows directly from \cref{lem:seed}, and isolates the
\emph{contrapositive} fork between non-additive seeds and nonzero
Leibniz defect as a separate remark.

\begin{proposition}[Diagnostic; corollary of \cref{lem:seed}]\label{prop:diagnostic}
Let $\mathcal N$ be a black noise and let
$(\delta_{\mathcal N},D_{\mathcal N})$ be an energy lift over
$\mathcal N$ in the sense of \cref{def:energy-lift}. Suppose
$\delta_{\mathcal N}$ has a deterministic-coefficient seed
\[
  T_{s,t}(h)\;:=\;\delta_{s,t}(\mathbf 1\otimes h),\qquad h\in H_{s,t},
\]
that is interval-additive in the sense of \cref{def:seed}. Then the
centered seed is trivial in $L^2$:
\begin{equation}\label{eq:trivial-seed}
  T_{s,t}(h_{s,t})\;=\;\E[T_{s,t}(h_{s,t})]\qquad\text{for every compatible }h.
\end{equation}
\end{proposition}

\begin{proof}
Direct from \cref{lem:seed}: the family $(T_{s,t})$ satisfies the
hypotheses of the lemma, so the centered family $X_{s,t}$ of
\eqref{eq:X-family} lies in $I_1(\mathcal N)$. Blackness of
$\mathcal N$ makes $I_1(\mathcal N)$ trivial, hence
$X_{s,t}=0$ in $L^2$, which is \eqref{eq:trivial-seed}.
\end{proof}

\begin{remark}[The fork]\label{rem:diagnostic-fork}
\Cref{prop:diagnostic} kills the interval-additive
deterministic-coefficient seed in black noise. It does not, by
itself, address whether $D_{\mathcal N}$ is a derivation. The
following two pressure points are therefore complementary:
\begin{itemize}
\item[$(\mathrm A)$] \emph{Seed pressure point.} For the calculus
to have nontrivial stochastic content via a deterministic-coefficient
seed, the seed must not be interval-additive in the sense of
\cref{def:seed}. It must be non-additive in $(s,t)$, nonlinear,
multiplicative, quadratic, or of area- or bracket-type, in the
senses left open by \cref{rem:seed-interpretation}.
\item[$(\mathrm B)$] \emph{Leibniz pressure point.} For the calculus
to have a nontrivial algebraic structure on a test algebra, the
adjoint derivative $D_{\mathcal N}$ either is a derivation
($\mathfrak L_{D_{\mathcal N}}\equiv 0$) or has a nonzero Leibniz
defect. The Poisson precedent \eqref{eq:leibniz-defect-poisson}
shows that nonzero defects are the algebraic signature of
non-first-order calculi already inside the classical L\'evy world,
and are the natural location for a nonclassical replacement of the
Wiener Leibniz rule.
\end{itemize}
Any nontrivial energy-first calculus over a black noise must engage
$(\mathrm A)$ nontrivially. Whether it must also engage $(\mathrm B)$
nontrivially --- that is, whether blackness forces a nonzero Leibniz
defect on any separating test algebra --- is open. The present paper
does not assert any necessity here; it asserts only that the
defect is the natural diagnostic.
\end{remark}

\begin{remark}[Diagnostic, not impossibility]\label{rem:leibniz-diagnostic}
\Cref{prop:diagnostic} is a diagnostic. It does not claim that no
$(\delta_{\mathcal N},D_{\mathcal N})$ exists over a black noise; it
does not claim that blackness forces
$\mathfrak L_{D_{\mathcal N}}\not\equiv 0$. It claims only that the
classical pattern --- interval-additive deterministic seed --- has
no nontrivial instance over a black noise, and that a nonzero
Leibniz defect is the natural place to look for a non-classical
replacement, on the basis of the Poisson precedent
\eqref{eq:leibniz-defect-poisson}.
\end{remark}

\subsection{\texorpdfstring{Carr\'e-du-champ as a positive seed}{Carre-du-champ as a positive seed}}\label{sec:cdc-seed}

The seed lemma (\cref{lem:seed}) rules out centered interval-additive
deterministic seeds over a black noise. It does not rule out
second-order energy seeds. The natural positive replacement is the
\emph{carr\'e-du-champ} bilinear form of Bakry--\'Emery, which we
record here as the principal candidate primitive datum for an
operator-derivative lift over a nonclassical noise.

For a generator $L$ on a function algebra $\mathcal A$ closed under
products, the carr\'e-du-champ is the bilinear form
\begin{equation}\label{eq:carre-du-champ}
  \Gamma(F,G)\;:=\;\tfrac{1}{2}\bigl(L(FG)-F\,LG-G\,LF\bigr),
  \qquad F,G\in\mathcal A.
\end{equation}
The form is symmetric in $(F,G)$, vanishes when either argument is a
constant, and is positive semidefinite whenever $L$ is the generator
of a Markov semigroup on a suitable algebra
(see~\cite{BakryGentilLedoux}, Ch.~1).

\paragraph{Why $\Gamma$ is not ruled out by the seed lemma.}
\Cref{lem:seed} concerns \emph{first-order} interval-additive seeds:
families $T_{s,t}\colon H_{s,t}\to L^{2}(\mathcal F_{s,t})$ linear in
the test direction $h\in H_{s,t}$, additive under
$H_{r,t}=H_{r,s}\oplus H_{s,t}$, and centered. The carr\'e-du-champ
is intrinsically \emph{second-order} in the test arguments: it is a
bilinear function of the pair $(F,G)$, with no requirement of
interval additivity in either slot. It specifies a positive energy
kernel rather than a first-chaos increment. Over a black noise the
additive first-chaos input vanishes by \cref{lem:seed}, but
$\Gamma$, as a bilinear positive form on a test algebra, need not.

\paragraph{Carr\'e-du-champ as a positive kernel.}
For a wide class of nonclassical generators, $\Gamma$ admits a
positive-kernel representation. For the fractional Laplacian
$L=-(-\Delta)^{\alpha/2}$ on $\R^d$, $\alpha\in(0,2)$, the
carr\'e-du-champ at a smooth $F$ takes the nonlocal form
\begin{equation}\label{eq:cdc-frac}
  \Gamma(F,F)(x)
  \;=\;c_{d,\alpha}\int_{\R^{d}}\frac{(F(x+h)-F(x))^{2}}{|h|^{d+\alpha}}\,dh,
\end{equation}
which is manifestly nonnegative. The integral kernel
$|h|^{-d-\alpha}$ can be interpreted in Fourier variables through a
positive covariance-type kernel; a separate treatment of this
identification in the energy-first framework appears in the
companion work~\cite{FontesFracBE}. The point for the present paper
is conservative: $\Gamma$ is not merely the algebraic remainder of a
Leibniz failure but a positive kernel encoding an intrinsic energy
structure of $L$, and this energy structure persists in regimes where
the classical first-chaos seed is unavailable.

\paragraph{Two-point analogue.}
For a stochastic flow of kernels $\Phi$ with one- and two-point
generators $L^{(1)}$, $L^{(2)}$, the natural two-point
carr\'e-du-champ on product observables is
\begin{equation}\label{eq:gamma-twopoint}
  \Gamma^{(2)}(f,g)(x,y)\;:=\;\tfrac{1}{2}\,\mathfrak D^{(2)}(f,g)(x,y),
\end{equation}
where $\mathfrak D^{(2)}$ is the two-point product defect of
\cref{def:two-point-defect} (\cref{sec:twopoint}). The factor
$\tfrac{1}{2}$ aligns the normalization with Bakry--\'Emery: in
coalescing-type examples, when the two-point dynamics reduce on the
diagonal to the one-point dynamics, $\Gamma^{(2)}(f,g)(x,x)$
coincides with the one-point carr\'e-du-champ $\Gamma(f,g)(x)$ built
from $L^{(1)}$, as in \eqref{eq:defect-coalescing-diag}. In examples
where the two-point dynamics do not reduce on the diagonal --- such
as common-driving flows (\cref{prop:defect-common}), where
$\Gamma^{(2)}$ is supported on all of $\R^{2}$ --- the diagonal
restriction is not pinned by the one-point semigroup. The product
defect of \cref{sec:twopoint} is in either case precisely a
two-point carr\'e-du-champ in the normalization
\eqref{eq:gamma-twopoint}; the conditional sticky-flow proposition
(\cref{prop:defect-sticky}) is a conditional identification of
$\Gamma^{(2)}_{\theta}$ as a diagonal-supported product-defect
distribution.

\paragraph{The primary candidate primitive datum.}
The seed taxonomy of \cref{sec:seed-taxonomy} below lists four
candidate seed classes. Among them, the carr\'e-du-champ /
product-defect class is the one for which explicit positive formulas
are available in nonclassical settings: the fractional-Laplacian
kernel \eqref{eq:cdc-frac} for nonlocal generators, the two-point
carr\'e-du-champ \eqref{eq:gamma-twopoint} for Le Jan--Raimond
flows. The multiplicative, area/Chen-type, and Boolean-projection
classes are at this stage less concrete. We therefore promote the
carr\'e-du-champ / product-defect class to \emph{primary} status in
the operator-derivative program over a nonclassical noise: the
candidate primitive datum is a positive bilinear energy kernel on a
test algebra, and the closed adjoint pair
$(\delta_{\mathcal N},D_{\mathcal N})$ is to be sought
\emph{consistent with} that kernel rather than constructed from a
first-chaos generator.

%==============================================================
\section{A two-point product defect for sticky flows}\label{sec:twopoint}
%==============================================================

The two pressure points of \cref{sec:seed,sec:leibniz} identify what
a primitive divergence over a black noise cannot reduce to: it
cannot be reached from a centered interval-additive deterministic
seed, and its adjoint operator derivative cannot in general be
expected to be first-order on a separating test algebra. The
pressure points do not, on their own, indicate what positive seed
data might support a construction. This section records the first
concrete candidate: the semigroup-level product defect of a
stochastic flow of kernels, specialized to the Le Jan--Raimond
sticky and coalescing flows~\cite{LeJanRaimond2004}.

The object we compute is not the operator-derivative Leibniz defect
of \cref{def:leibniz-defect}, which would require a primitive
divergence $\delta_{\mathcal N}$ in hand. It is the
semigroup-level \emph{shadow} of that defect, accessible already at
the level of the one- and two-point semigroups of the flow. The
shadow is computable and structurally explicit on standard
examples, and it naturally singles out a class of nonclassical seed
data that \cref{lem:seed} does not rule out.

\subsection{The two-point product defect}\label{sec:twopoint-def}

Let $\Phi$ be a stochastic flow of kernels on $\R$ in the sense of
\cite{LeJanRaimond2004}, with associated noise $\mathcal N_{\Phi}$.
Let $\{P_t^{(1)}\}_{t\ge 0}$ denote its one-point semigroup acting
on a test class $\mathcal T^{(1)}\subset C_b(\R)$ with generator
$L^{(1)}$, and let $\{P_t^{(2)}\}_{t\ge 0}$ denote its two-point
semigroup acting on a test class $\mathcal T^{(2)}\subset C_b(\R^2)$
with generator $L^{(2)}$. We assume product observables
$f\otimes g(x,y)=f(x)g(y)$ with $f,g\in\mathcal T^{(1)}$ lie in
$\mathcal T^{(2)}$.

\begin{definition}[Two-point product defect]\label{def:two-point-defect}
The \emph{two-point product defect} of the flow $\Phi$ on
$f,g\in\mathcal T^{(1)}$ is the function
\begin{equation}\label{eq:two-point-defect}
  \mathfrak D^{(2)}(f,g)(x,y)
  \;:=\;
  L^{(2)}(f\otimes g)(x,y)
  -(L^{(1)}f)(x)\,g(y)
  -f(x)\,(L^{(1)}g)(y),
\end{equation}
read on $\R^{2}$ (or on its symmetric quotient as appropriate).
\end{definition}

\Cref{def:two-point-defect} measures the failure of $L^{(2)}$ to
act as the tensor product
$L^{(1)}\otimes\mathrm{id}+\mathrm{id}\otimes L^{(1)}$ on product
observables. It is the semigroup-level analogue, on product
observables built from one-point evaluations, of the
operator-derivative Leibniz defect $\mathfrak L_{D_{\mathcal N}}$
of \cref{def:leibniz-defect}. When the noise of the flow is
classical and a Malliavin pair $(\delta^{\Phi},D^{\Phi})$ is
available, $\mathfrak D^{(2)}$ is recovered from
$\mathfrak L_{D^{\Phi}}$ on product cylindrical functionals by the
standard generator-to-Malliavin translation; over a black noise the
operator-derivative Leibniz defect is not constructed, but the
semigroup-level shadow $\mathfrak D^{(2)}$ remains a computable
invariant of the flow.

\subsection{Independent Brownian motion benchmark}\label{sec:twopoint-indep}

The simplest benchmark is the flow of independent Brownian motions,
in which each starting point evolves as an independent Brownian
motion and the two-point semigroup is the product of one-point
semigroups.

\begin{proposition}[Zero defect for independent Brownian motion]\label{prop:defect-indep}
Let $\Phi$ be the stochastic flow of independent Brownian motions
on $\R$, with $L^{(1)}=\tfrac{1}{2}\partial_x^2$ and
$L^{(2)}=\tfrac{1}{2}(\partial_x^2+\partial_y^2)$ on
$\mathcal T^{(1)}=\mathcal T^{(2)}=C_c^{\infty}$. Then
\begin{equation}\label{eq:defect-indep-zero}
  \mathfrak D^{(2)}(f,g)(x,y)\;\equiv\;0,
  \qquad f,g\in C_c^{\infty}(\R),\ (x,y)\in\R^2.
\end{equation}
\end{proposition}

\begin{proof}
$L^{(2)}(f\otimes g)(x,y)=\tfrac{1}{2}f''(x)g(y)+\tfrac{1}{2}f(x)g''(y)
=(L^{(1)}f)(x)g(y)+f(x)(L^{(1)}g)(y)$.
\end{proof}

\subsection{Common-driving and coalescing benchmarks}\label{sec:twopoint-coalescing}

A flow with off-diagonal interaction produces a nonzero defect. Two
standard cases place the interaction in qualitatively different
locations: globally for a common-driving flow, and on the diagonal
only for a coalescing flow.

\begin{proposition}[Common-driving flow: global defect]\label{prop:defect-common}
Let $\Phi$ be the flow on $\R$ driven by a single common Brownian
motion $W$, $X_t^{(x)}=x+W_t$, with two-point generator
$L^{(2)}=\tfrac{1}{2}(\partial_x+\partial_y)^2
=\tfrac{1}{2}(\partial_x^2+2\partial_{xy}+\partial_y^2)$. Then
\begin{equation}\label{eq:defect-common}
  \mathfrak D^{(2)}(f,g)(x,y)\;=\;f'(x)\,g'(y),\qquad (x,y)\in\R^2.
\end{equation}
\end{proposition}

\begin{proof}
$L^{(2)}(f\otimes g)(x,y)=\tfrac{1}{2}f''(x)g(y)+f'(x)g'(y)+\tfrac{1}{2}f(x)g''(y)$;
subtracting the one-point contributions leaves $f'(x)g'(y)$.
\end{proof}

\begin{proposition}[Coalescing Brownian flow: formal diagonal defect]\label{prop:defect-coalescing}
Let $\Phi$ be Arratia's coalescing Brownian flow on $\R$, in which
distinct points evolve as independent Brownian motions until they
meet and coalesce, and as a single Brownian motion thereafter. For
$f,g\in C_c^{\infty}(\R)$:
\begin{itemize}
\item[$(\mathrm a)$] \emph{Off-diagonal vanishing.} On
$\{(x,y)\in\R^{2}:x\neq y\}$,
\begin{equation}\label{eq:defect-coalescing-offdiag}
  \mathfrak D^{(2)}(f,g)(x,y)\;=\;0.
\end{equation}
\item[$(\mathrm b)$] \emph{Diagonal density.} On the diagonal
$\{x=y\}$, with the identification
$(f\otimes g)(x,x)=(fg)(x)$ inherited from the coalescence rule,
the two-point generator acts as
$L^{(2)}(f\otimes g)(x,x)=L^{(1)}(fg)(x)=\tfrac{1}{2}(fg)''(x)$, so
the pointwise diagonal value of the product defect is the one-point
carr\'e-du-champ of the heat semigroup,
\begin{equation}\label{eq:defect-coalescing-diag}
  \mathfrak D^{(2)}(f,g)(x,x)
  \;=\;f'(x)g'(x)
  \;=\;2\,\Gamma^{(1)}(f,g)(x).
\end{equation}
\end{itemize}
In particular $\mathfrak D^{(2)}(f,g)$ is a product-defect
distribution supported on the diagonal $\{x=y\}$, with formal density
$f'(x)g'(x)$ in the variable $x$. The normalization of
$\mathfrak D^{(2)}(f,g)$ as a distribution on $\R^{2}$ --- against the
diagonal Lebesgue measure, against the local-time measure of
$Y_t-X_t$ at $0$, or against any other choice --- is inherited from
the generator and domain convention adopted for the coalescing flow.
\end{proposition}

\begin{proof}
For $(\mathrm a)$, off the diagonal the two distinct points of the
coalescing flow evolve as independent Brownian motions and the
defect vanishes by \cref{prop:defect-indep}. For $(\mathrm b)$, the
coalescence identification gives
$L^{(2)}(f\otimes g)(x,x)=L^{(1)}(fg)(x)=\tfrac{1}{2}(fg)''(x)$,
while $(L^{(1)}f)(x)g(x)+f(x)(L^{(1)}g)(x)
=\tfrac{1}{2}f''(x)g(x)+\tfrac{1}{2}f(x)g''(x)$, and subtraction
yields $f'(x)g'(x)$. The relation with the one-point carr\'e-du-champ
$\Gamma^{(1)}(f,g)(x)=\tfrac{1}{2}(L^{(1)}(fg)-fL^{(1)}g-gL^{(1)}f)(x)=\tfrac{1}{2}f'(x)g'(x)$
of \eqref{eq:carre-du-champ} is then $\mathfrak D^{(2)}(f,g)(x,x)=2\Gamma^{(1)}(f,g)(x)$.
\end{proof}

The two benchmarks illustrate the qualitative spectrum of
$\mathfrak D^{(2)}$. For independent flows it vanishes identically.
For a common-driving flow it is a globally nonzero bilinear form
supported on $\R^2$. For a coalescing flow it is concentrated on
the measure-zero diagonal, recording that the two-point interaction
is confined to the contact set. The Brownian sticky flow sits
between the independent and coalescing regimes on the diagonal.

\subsection{Conditional proposition for the Brownian sticky flow}\label{sec:twopoint-sticky}

The Brownian sticky flow constructed in \cite{LeJanRaimond2004} is
a one-parameter interpolation $\Phi_{\theta}$, with one endpoint of
the parametrization recovering the flow of independent Brownian
motions and the other recovering the coalescing Brownian flow, and
intermediate values producing sticky encounters at the diagonal.
The off-diagonal two-point dynamics are independent in both Brownian
coordinates; the sticky structure lies entirely at the boundary
$\{x=y\}$. The product defect $\mathfrak D^{(2)}_{\theta}$ inherits
this structure.

\begin{proposition}[Brownian sticky flow: conditional defect]\label{prop:defect-sticky}
Let $\Phi_{\theta}$ be the Brownian sticky flow of
\cite{LeJanRaimond2004} with sticky parameter $\theta$. Suppose
the two-point generator $L^{(2)}_{\theta}$ acts on a suitable test
class as $\tfrac{1}{2}(\partial_x^2+\partial_y^2)$ off the diagonal
$\{x=y\}$, subject to the Wentzell-type sticky boundary condition
prescribed by \cite{LeJanRaimond2004}. Then the product defect of
\cref{def:two-point-defect} vanishes off the diagonal,
\begin{equation}\label{eq:defect-sticky-offdiag}
  \mathfrak D^{(2)}_{\theta}(f,g)(x,y)\;=\;0\qquad\text{for }x\neq y,
\end{equation}
and is concentrated as a diagonal-supported product-defect
distribution on $\{x=y\}$. The diagonal coefficient is governed by
the sticky boundary condition; it vanishes at the independent-flow
endpoint of the parametrization (recovering
\cref{prop:defect-indep}), and at the coalescing endpoint of the
parametrization the diagonal formal density is $f'(x)g'(x)$ in the
sense of \cref{prop:defect-coalescing}, with normalization
determined by the chosen generator/domain convention.
\end{proposition}

The off-diagonal statement \eqref{eq:defect-sticky-offdiag} follows
from \cref{prop:defect-indep} applied to the off-diagonal dynamics
of $\Phi_{\theta}$. The diagonal statement is conditional on the
exact form of the Wentzell-type boundary condition in
\cite{LeJanRaimond2004}; we record this as the verification step.

\begin{remark}[Verification requirement]\label{rem:lejanraimond-todo}
\Cref{prop:defect-sticky} is stated conditionally on the explicit
two-point sticky generator $L^{(2)}_{\theta}$ and on its
Wentzell-type boundary condition at $\{x=y\}$. The data required to
make the diagonal coefficient explicit, and hence
\cref{prop:defect-sticky} unconditional, are: (a) the domain of
$L^{(2)}_{\theta}$ as an unbounded operator together with the
boundary condition characterizing it; (b) the parametrization
convention for $\theta$ and the identification of the independent
and coalescing limits; (c) the precise rescaling of the diagonal
occupation measure (a local time of the difference $Y_t-X_t$ at
$0$) and the explicit $\theta$-dependence of the diagonal
coefficient. We refer to \cite{LeJanRaimond2004} for these data and
treat their extraction as the first concrete verification step for
the present approach.
\end{remark}

\subsection{Positive seed taxonomy}\label{sec:seed-taxonomy}

The propositions of \S\ref{sec:twopoint-coalescing}--\ref{sec:twopoint-sticky}
identify a class of positive seed candidates that \cref{lem:seed}
does not rule out: the \emph{carr\'e-du-champ / product-defect
seeds} of \cref{sec:cdc-seed}. We record this class as the primary
candidate, together with three further candidate types, as a
taxonomy of seed data for a primitive divergence over a nonclassical
noise. The list is descriptive: none of these candidates has been
extended to a closed $\delta_{\mathcal N}$ over a nontrivial black
noise; the point of listing them is to fix the vocabulary in which
partial answers to \cref{prog:open} can be classified.

\paragraph{Carr\'e-du-champ / product-defect seeds (primary).}
A symmetric bilinear form
$\Gamma_{\mathcal N}\colon\mathcal A\times\mathcal A\to L^{2}(\mathcal N)$
on a test algebra $\mathcal A\subset L^{\infty}(\mathcal N)$,
specified intrinsically from the noise tensor factorization in the
manner of \eqref{eq:carre-du-champ}, with the energy-form positivity
of \cref{prog:open}$(\mathrm i)$. The concrete instances are
\eqref{eq:gamma-twopoint} for stochastic flows of kernels --- the
two-point carr\'e-du-champ extracted from
\cref{prop:defect-indep,prop:defect-common,prop:defect-coalescing,prop:defect-sticky}
--- and \eqref{eq:cdc-frac} for nonlocal generators. This is the
only seed class for which explicit positive formulas are available
in nonclassical settings.

\paragraph{Multiplicative seeds.}
A family $(M_{s,t})$ of $\mathcal F_{s,t}$-measurable elements of
$L^{2}(\mathcal N)$ satisfying the multiplicative interval law
\[
  M_{r,t}\;=\;M_{r,s}\,M_{s,t},\qquad r\le s\le t,
\]
rather than the additive law of \eqref{eq:additivity}. The classical
analogue is the stochastic exponential, which factors through a
logarithmic first chaos; a multiplicative seed is not required to
do so, and so is not trivialized by \cref{lem:seed}.

\paragraph{Area- or Chen-type seeds.}
A second-level family $(A_{r,t})$ satisfying a Chen-type relation
$A_{r,t}=A_{r,s}+A_{s,t}+B_{r,s}\otimes B_{s,t}$ together with a
(possibly trivial) first-level $(B_{s,t})$. Over a black noise the
first level is forced to be trivial by \cref{lem:seed}; the
second-level $A_{r,t}$ is then either purely additive or carries an
intrinsically defined second-level interaction independent of any
first level.

\paragraph{Boolean-projection seeds.}
A seed defined through the family of $L^{2}$-conditional-expectation
projections $\{Q_{\mathcal H}\}_{\mathcal H\in\mathfrak B(\mathcal N)}$
associated with the noise Boolean algebra of $\mathcal N$, rather
than through real-valued interval families. Over a black noise the
projection-differential first chaos is trivial; a Boolean-projection
seed must come from higher-order algebraic combinations on the
Boolean lattice. We refer to \cite{Tsirelson2004,TsirelsonVershik1998}
for the Boolean formulation of noise structure and treat the
Boolean-projection seed as a coordinate-system proposal rather than
a constructed object.

\medskip
The carr\'e-du-champ / product-defect class is the most concrete
member of this taxonomy and the natural target for the first
verification step of \cref{prog:open}. The remaining three classes
are listed for completeness; their extension to constructed calculi
over black noises is open and not addressed in the present paper.

\subsection{Relation to the operator-derivative Leibniz defect}\label{sec:twopoint-relation}

The two-point product defect $\mathfrak D^{(2)}$ and the
operator-derivative Leibniz defect $\mathfrak L_{D_{\mathcal N}}$
are objects of different levels of structure. The former is a
semigroup-level invariant of a flow, computable from $L^{(1)}$ and
$L^{(2)}$ on product observables. The latter is the algebraic
invariant of a primitive divergence's adjoint, requiring
$(\delta_{\mathcal N},D_{\mathcal N})$ in hand. We make the relation
between them explicit.

\begin{remark}[\texorpdfstring{$\mathfrak D^{(2)}$ as the shadow of $\mathfrak L_{D_{\mathcal N}}$}{D^(2) as the shadow of L_D}]\label{rem:defect-shadow}
Suppose, for the duration of this remark only, that an
operator-derivative lift
$(\delta_{\mathcal N_{\Phi}},D_{\mathcal N_{\Phi}})$ of the flow
$\Phi$ exists, with $D_{\mathcal N_{\Phi}}$ acting on a test algebra
$\mathcal A_{\Phi}$ containing the product cylindrical functionals
$F_{f,g}^{(x,y)}:=f(X_t^{(x)})g(X_t^{(y)})$ for $f,g\in\mathcal T^{(1)}$
and $x,y\in\R$. Then on $F_{f,g}^{(x,y)}$ the operator-derivative
Leibniz defect $\mathfrak L_{D_{\mathcal N_{\Phi}}}$ and the
semigroup-level defect $\mathfrak D^{(2)}$ are connected by the
standard generator-to-Malliavin translation: the bilinear coefficient
of $\mathfrak D^{(2)}(f,g)$ recovers, after this translation, the
$L^{\infty}(\mathcal N_{\Phi})$-bilinear component of
$\mathfrak L_{D_{\mathcal N_{\Phi}}}$ on the product algebra
generated by $F_f^{(x)}$. The construction of
$(\delta_{\mathcal N_{\Phi}},D_{\mathcal N_{\Phi}})$ that would make
this translation rigorous is the content of \cref{prog:open}; until
it is in hand, $\mathfrak D^{(2)}$ is the semigroup-level shadow
that any constructed $\mathfrak L_{D_{\mathcal N_{\Phi}}}$ should
recover or refine.
\end{remark}

\Cref{rem:defect-shadow} formulates the first concrete adequacy
criterion for the operator-derivative program on a Le Jan--Raimond
flow: any candidate
$(\delta_{\mathcal N_{\Phi}},D_{\mathcal N_{\Phi}})$ must recover
the explicit defects of
\cref{prop:defect-indep,prop:defect-common,prop:defect-coalescing,prop:defect-sticky}
on product observables. This criterion is independent of any first-chaos
data on $\mathcal N_{\Phi}$ and is therefore a meaningful target
even when, as for the coalescing and sticky flows, the noise of the
flow is black.

\medskip
\noindent\emph{Summary.} The product defect $\mathfrak D^{(2)}(f,g)$
of \cref{def:two-point-defect} is the two-point carr\'e-du-champ of
the flow on product observables, up to the Bakry--\'Emery
normalization \eqref{eq:gamma-twopoint}:
$\Gamma^{(2)}(f,g)=\tfrac{1}{2}\mathfrak D^{(2)}(f,g)$. For
independent Brownian motions it vanishes; for the common-driving
flow it is a globally nonzero positive kernel of formal density
$\tfrac{1}{2}f'(x)g'(y)$ on $\R^{2}$; for the coalescing flow it is
a diagonal-supported product-defect distribution of formal density
$\tfrac{1}{2}f'(x)g'(x)$ in the diagonal variable, with
normalization determined by the two-point generator and domain
convention (\cref{prop:defect-coalescing}); for the Brownian sticky
flow it is, conditionally on \cite{LeJanRaimond2004}, a
diagonal-supported product-defect distribution encoded by the
sticky boundary condition (\cref{prop:defect-sticky}). In each case
the candidate seed datum for the operator-derivative program over
the flow is not an additive first-chaos increment but a two-point
carr\'e-du-champ $\Gamma^{(2)}$ extracted from the two-point
semigroup. The seed lemma rules out the additive route; the
carr\'e-du-champ supplies the next-most-concrete positive candidate.

\subsection{\texorpdfstring{From two-point energy to the \(N\)-point defect hierarchy}{From two-point energy to the N-point defect hierarchy}}\label{sec:twopoint-hierarchy}

The two-point product defect of \cref{def:two-point-defect} is one
layer of a more general construction. Whenever a flow or noise
carries $N$-point semigroups $\{P_t^{(N)}\}_{t\ge 0}$ with generators
$L^{(N)}$ acting on a test class
$\mathcal T^{(N)}\subset C_b(\R^N)$ that contains the product
cylinders $f_1\otimes\cdots\otimes f_N$ with
$f_i\in\mathcal T^{(1)}$, we record the natural higher analogue.

\begin{definition}[\texorpdfstring{\(N\)-point product defect}{N-point product defect}]\label{def:N-point-defect}
The \emph{$N$-point product defect} of the flow $\Phi$ on
$f_1,\ldots,f_N\in\mathcal T^{(1)}$ is the object
\begin{equation}\label{eq:N-point-defect}
  \mathfrak D^{(N)}(f_1,\ldots,f_N)
  \;:=\;
  L^{(N)}(f_1\otimes\cdots\otimes f_N)
  \;-\;\sum_{i=1}^{N}(L^{(1)}f_i)\otimes\prod_{j\neq i}f_j,
\end{equation}
initially defined on the common domain of the $N$-point generator
$L^{(N)}$ and of the one-point generator $L^{(1)}$ acting in each
factor, and interpreted according to the $N$-point generator/domain
convention --- as a function, an energy measure, or a distribution
on $\R^{N}$ (or on its symmetric quotient as appropriate). In
particular the diagonal-supported product defects of the coalescing
and sticky flows in
\cref{prop:defect-coalescing,prop:defect-sticky} are read as
distributions on $\R^{2}$.
\end{definition}

For $N=2$ this recovers \cref{def:two-point-defect}, the failure of
$L^{(2)}$ to equal
$L^{(1)}\otimes\mathrm{id}+\mathrm{id}\otimes L^{(1)}$. For $N\ge 3$
the raw defect $\mathfrak D^{(N)}$ already contains contributions
from proper subsets $B\subsetneq\{1,\ldots,N\}$ lifted to all $N$
variables. To isolate the \emph{genuinely $N$-body} interaction we
work at the infinitesimal generator level: we organize
$\mathfrak D^{(N)}$ over subsets of $\{1,\ldots,N\}$ on the Boolean
lattice of subsets, rather than through a finite-time partition
expansion. The right index for the hierarchy is the finite subset
$S$ itself, not the integer $|S|$.

\subsubsection*{Subset-indexed half raw product defect}

For a finite index set $S$ with $|S|=N\ge 2$ and a tuple of one-point
test functions $(f_i)_{i\in S}\subset\mathcal T^{(1)}$, write
$f_S:=\prod_{i\in S}f_i$ for the product observable and $L^{(S)}$
for the $|S|$-point generator of the flow acting on $f_S$. Define the
\emph{half raw product defect} on $S$ by
\begin{equation}\label{eq:M-S}
  M_S(f_S)
  \;:=\;
  \frac{1}{2}\!\left[
    L^{(S)}(f_S)
    \;-\;
    \sum_{i\in S}(L^{(1)}f_i)\prod_{j\in S\setminus\{i\}}f_j
  \right]
  \;=\;\tfrac{1}{2}\mathfrak D^{(|S|)}(f_S),
\end{equation}
where the second equality matches the integer-indexed
\cref{def:N-point-defect} under the relabeling
$\{1,\ldots,|S|\}\leftrightarrow S$. For $|S|=2$,
$M_S(f_S)=\Gamma^{(2)}(f_i,f_j)$, the bilinear two-point
carr\'e-du-champ of \eqref{eq:gamma-twopoint}.

\begin{definition}[Generator-level interaction hierarchy]\label{def:N-point-cdc}
For a finite index set $S$ with $|S|\ge 2$ and one-point test
functions $(f_i)_{i\in S}\subset\mathcal T^{(1)}$, define the
\emph{generator-level interaction defect} $\Gamma_S$ on $S$
recursively by
\begin{equation}\label{eq:gamma-N-cumulant}
  \Gamma_S(f_S)
  \;:=\;
  M_S(f_S)
  \;-\;
  \!\!\sum_{\substack{B\subsetneq S\\ |B|\ge 2}}\!\!
  \Gamma_B(f_B)\prod_{j\in S\setminus B}f_j,
\end{equation}
read in the appropriate distribution / energy-measure / function
space on $\R^{|S|}$ inherited from the $|S|$-point generator/domain
convention. The base case $|S|=2$ gives
$\Gamma_{\{i,j\}}(f_if_j)=M_{\{i,j\}}(f_if_j)=\Gamma^{(2)}(f_i,f_j)$,
the bilinear two-point carr\'e-du-champ.
Equivalently, on every $S$ the recursion is the subset inversion of
\begin{equation}\label{eq:gamma-N-inversion}
  M_S(f_S)
  \;=\;
  \!\!\sum_{\substack{B\subseteq S\\ |B|\ge 2}}\!\!
  \Gamma_B(f_B)\prod_{j\in S\setminus B}f_j.
\end{equation}
We refer to $\{\Gamma_S\}_{|S|\ge 2}$ as the \emph{generator-level
interaction hierarchy} (or, equivalently, the
\emph{infinitesimal product-defect hierarchy}); the members
$\Gamma_S$ for $|S|\ge 3$ are the \emph{higher
$\Gamma$-interaction defects} of the flow. The bilinear $|S|=2$
member is the proper carr\'e-du-champ; the higher members are
$|S|$-linear interaction defects, not bilinear carr\'es-du-champ.
\end{definition}

The inversion \eqref{eq:gamma-N-inversion} sums over \emph{subsets}
of $S$, not over set partitions of $S$. The distinction is essential
and is the reason the hierarchy is described as generator-level
rather than as a cumulant family.

\begin{remark}[Partition cumulants vs.\ generator-level defects]\label{rem:partition-vs-generator}
In finite-time moment and cumulant expansions one must sum over set
partitions of $S$. At $|S|=4$ the partition inversion contains
terms of partition-product type such as
$\Gamma^{(2)}_{\{1,2\}}\,\Gamma^{(2)}_{\{3,4\}}$ ---
products of two lower interactions arising from $t^{2}$ and higher
contributions in $P_t^{(N)}=e^{tL^{(N)}}$ acting on a product
observable. Such terms belong to finite-time cumulants and to
higher-order-in-time expansions; they are bilinear in the
interaction hierarchy.
The present paper instead decomposes the \emph{infinitesimal}
product defect --- the first-order $dt$-content of $L^{(N)}$ on
$\prod_{i\in S}f_i$ --- and at this generator level the
decomposition is linear over interacting subsets of $S$, not
bilinear over set partitions of $S$. The recursion
\eqref{eq:gamma-N-cumulant} is correspondingly a subset inversion,
and no partition-product term appears. A finite-time
partition-cumulant hierarchy may be derived from the semigroups
$\{P_t^{(N)}\}_{t\ge 0}$ at a later stage; it is not the object of
this paper.
\end{remark}

\paragraph{Schematic for small $|S|$.}
For $|S|=2$, $\Gamma_{\{i,j\}}(f_if_j)=M_{\{i,j\}}(f_if_j)$. For
$S=\{1,2,3\}$,
\begin{equation}\label{eq:gamma-3}
  \Gamma_{\{1,2,3\}}(f_1f_2f_3)
  \;=\;
  M_{\{1,2,3\}}(f_1f_2f_3)
  -\Gamma_{\{1,2\}}(f_1f_2)\,f_3
  -\Gamma_{\{1,3\}}(f_1f_3)\,f_2
  -\Gamma_{\{2,3\}}(f_2f_3)\,f_1,
\end{equation}
in which the three pairwise interaction defects are subtracted from
$M_{\{1,2,3\}}$ to leave the irreducible three-body content. For
$S=\{1,2,3,4\}$ (writing $\Gamma_B$ for $\Gamma_B(f_B)$ to lighten
notation),
\begin{equation}\label{eq:gamma-4}
\begin{aligned}
\Gamma_{\{1,2,3,4\}}
&=
M_{\{1,2,3,4\}}\\
&\quad
-\Gamma_{\{1,2,3\}}f_4-\Gamma_{\{1,2,4\}}f_3-\Gamma_{\{1,3,4\}}f_2-\Gamma_{\{2,3,4\}}f_1\\
&\quad
-\Gamma_{\{1,2\}}f_3f_4-\Gamma_{\{1,3\}}f_2f_4-\Gamma_{\{1,4\}}f_2f_3\\
&\quad
-\Gamma_{\{2,3\}}f_1f_4-\Gamma_{\{2,4\}}f_1f_3-\Gamma_{\{3,4\}}f_1f_2,
\end{aligned}
\end{equation}
in which the ten subtracted terms enumerate the four three-body
and six two-body interaction defects on the proper subsets
$B\subsetneq\{1,2,3,4\}$ with $|B|\ge 2$. \emph{No
partition-product term such as
$\Gamma_{\{1,2\}}\Gamma_{\{3,4\}}$ appears in the generator-level
hierarchy}; such terms belong to finite-time partition cumulants
and not to the infinitesimal product defect, in the sense of
\cref{rem:partition-vs-generator}.

\paragraph{Distribution-valued caution.}
The $\Gamma_S$ are interpreted in the appropriate space inherited
from the $|S|$-point generator/domain convention. In regular cases
they are represented by functions on $\R^{|S|}$; in coalescing or
sticky cases they may be diagonal-supported measures or
distributions. The same caution applies to the half raw defects
$M_S$.

\paragraph{Classical truncation.}
For a second-order diffusion on $\R^{d}$ whose $|S|$-point
dynamics are generated by sums of one- and two-particle terms ---
products of independent copies, common drivers, finite-range
pairwise couplings, or coalescence rules that act only at pair
encounters --- the half raw defect satisfies the pair-generated
identity
\begin{equation}\label{eq:pair-generated}
  M_S(f_S)
  \;=\;
  \!\!\sum_{\substack{B\subseteq S\\ |B|=2}}\!\!
  \Gamma_B(f_B)\prod_{j\in S\setminus B}f_j,
  \qquad |S|\ge 2.
\end{equation}
Comparing with the full inversion \eqref{eq:gamma-N-inversion}
forces $\Gamma_S=0$ for every $|S|\ge 3$, up to the existing caveat
about higher collision strata and the diagonal/boundary conditions
specifying $L^{(|S|)}$. In that case the hierarchy truncates at the
two-point level. The benchmarks of
\cref{prop:defect-indep,prop:defect-common,prop:defect-coalescing}
sit in this class: their $|S|$-point generators are determined by
the one- and two-point laws, and the hierarchy is pair-saturated.

\paragraph{Nonclassical noises.}
There is no a priori reason for the hierarchy to truncate when the
$|S|$-point dynamics of a flow or noise are not pairwise-generated.
Stochastic flows with genuine multi-time, multi-particle interaction,
and noises constructed from extended random structures (multi-arm
events at criticality, multi-point coalescence with memory) can in
principle support nontrivial $\Gamma_S$ at every level $|S|\ge 2$.
We assert no such example here; the generator-level interaction
hierarchy is introduced as the natural all-$N$ extension of the
two-point seed, not as an existence statement on any particular
noise.

\paragraph{Hierarchy as candidate primitive datum.}
\Cref{sec:cdc-seed} promoted $\Gamma^{(2)}$ to the principal positive
candidate replacing the additive first-chaos seed. The natural
all-$|S|$ refinement is to take the full hierarchy
$\{\Gamma_{\mathcal N,S}\}_{|S|\ge 2}$ as candidate primitive
datum, with $\Gamma_{\mathcal N,\{i,j\}}$ --- the noise-level
two-point carr\'e-du-champ --- as its first and most concrete layer
(the layer for which the explicit positive formulas of
\cref{prop:defect-indep,prop:defect-common,prop:defect-coalescing,prop:defect-sticky}
and \eqref{eq:cdc-frac} are available) and the higher layers
$\Gamma_{\mathcal N,S}$ for $|S|\ge 3$ as the genuinely multi-body
generator-level algebraic content above pairwise energy. The open
problem of \cref{sec:open} is restated accordingly.

\begin{remark}[Non-claims for the hierarchy]\label{rem:hierarchy-non-claim}
We do not claim that the hierarchy
$\{\Gamma_{\mathcal N,S}\}_{|S|\ge 2}$ exists, intrinsically
specified from $\mathcal N$, on every black noise. We do not claim
that the hierarchy determines $\delta_{\mathcal N}$ even when it
exists, nor do we assert a representation theorem linking the
hierarchy to a closed adjoint pair
$(\delta_{\mathcal N},D_{\mathcal N})$. The two-point member
$\Gamma_{\mathcal N,\{i,j\}}$ remains the primary candidate seed
and the first concrete layer to be constructed in any partial answer
to \cref{prog:open}. A finite-time partition-cumulant hierarchy is
not constructed here either; see
\cref{rem:partition-vs-generator}.
\end{remark}

%==============================================================
\section{Open problem}\label{sec:open}
%==============================================================

We now state the central open problem of the paper, in the sharpened
form provided by the pressure-point analyses of
\cref{sec:seed,sec:leibniz} and the semigroup-level computation of
\cref{sec:twopoint}.

\begin{program}[\texorpdfstring{Primitive divergence and generator-level interaction hierarchy over a nonclassical noise}{Primitive divergence and generator-level interaction hierarchy over a nonclassical noise}]\label{prog:open}
Let $\mathcal N=(\mathcal F_{s,t})$ be a noise in the
Tsirelson--Vershik sense, possibly black. Specify a tuple
\[
  \bigl(\{\Gamma_{\mathcal N,S}\}_{|S|\ge 2},
        \delta_{\mathcal N},D_{\mathcal N}\bigr),
\]
indexed by finite subsets $S$, together with a test algebra
$\mathcal A\subset L^{\infty}(\mathcal N)$, satisfying:
\begin{itemize}
\item[$(\mathrm i)$] \emph{Primitive generator-level interaction
hierarchy.} A family of $|S|$-linear forms
\[
  \Gamma_{\mathcal N,S}\colon\mathcal A^{\,|S|}\longrightarrow L^{2}(\mathcal N),
  \qquad |S|\ge 2,
\]
indexed by finite subsets $S$, in which the bilinear member at
$|S|=2$ is the proper noise-level carr\'e-du-champ (denoted
$\Gamma_{\mathcal N}^{(2)}=\Gamma_{\mathcal N,\{i,j\}}$ when the
particular subset labels are immaterial) and, for $|S|\ge 3$,
$\Gamma_{\mathcal N,S}$ is the noise-level lift of the
generator-level interaction defect of \cref{def:N-point-cdc}. In
semigroup-level realizations the corresponding $|S|$-point objects
may be represented by functions, energy measures, or distributions;
the $L^{2}(\mathcal N)$-valued formulation refers to the intrinsic
noise-level object when such a realization exists.
The hierarchy is specified intrinsically from the tensor
factorization \eqref{eq:tensor-factorization} of $\mathcal N$,
without reference to a classical first-chaos generator. Each
$\Gamma_{\mathcal N,S}$ vanishes whenever any argument is a constant,
is compatible with the interval-tensor structure of $\mathcal N$,
and is required to be a generator-level interaction defect in the
subset-inversion sense of
\eqref{eq:gamma-N-cumulant}--\eqref{eq:gamma-N-inversion} with
respect to a generator hierarchy $L_{\mathcal N}^{(S)}$ associated
with $\mathcal N$ wherever that hierarchy is available. The
two-point form $\Gamma_{\mathcal N}^{(2)}$ is symmetric and
energy-form positive,
\[
  \E\!\bigl[\Gamma_{\mathcal N}^{(2)}(F,F)\bigr]\ge 0
  \quad\text{for every }F\in\mathcal A,
\]
with the pointwise sharpening
$\Gamma_{\mathcal N}^{(2)}(F,F)\ge 0$ almost surely when a pointwise
order is available on its values, and is identified as the primary
candidate seed. No positivity is required of $\Gamma_{\mathcal N,S}$
for $|S|\ge 3$. A finite-time partition-cumulant hierarchy may
later be derived from the semigroups; it is not part of this
formulation (see \cref{rem:partition-vs-generator}).
\item[$(\mathrm{ii})$] \emph{Primitive divergence.} A closed
densely-defined operator
$\delta_{\mathcal N}\colon\Dom(\delta_{\mathcal N})\subset L^{2}(\mathcal N;\mathcal U)\to L^{2}(\mathcal N)$,
compatible with the interval-tensor structure of $\mathcal N$ and
specified consistently with
$\{\Gamma_{\mathcal N,S}\}_{|S|\ge 2}$ in the sense of
$(\mathrm{iv})$ below.
\item[$(\mathrm{iii})$] \emph{Operator derivative.} The closed
adjoint $D_{\mathcal N}:=\delta_{\mathcal N}^{*}$ on the natural
dense domain in $L^{2}(\mathcal N)$, satisfying the
integration-by-parts identity \eqref{eq:IBP}. No nontrivial
classical first-chaos generator is used as an intermediate input.
\item[$(\mathrm{iv})$] \emph{Carr\'e-du-champ recovery, two-point
layer.} The formal Markov-type generator
$L_{\mathcal N}:=-\delta_{\mathcal N}D_{\mathcal N}$ built from the
pair is required to be densely defined on a common subdomain of
$\mathcal A$, and the Bakry--\'Emery identity
\[
  \Gamma_{\mathcal N}^{(2)}(F,G)
  \;=\;\tfrac{1}{2}\bigl(L_{\mathcal N}(FG)-F\,L_{\mathcal N}G-G\,L_{\mathcal N}F\bigr)
\]
is required to hold on $\mathcal A$. The Leibniz defect
$\mathfrak L_{D_{\mathcal N}}$ of \cref{def:leibniz-defect} on
$\mathcal A$ is identified explicitly as the algebraic remainder
relating the two sides. Recovery of the higher layers
$\Gamma_{\mathcal N,S}$ for $|S|\ge 3$ through a multi-point
generator-level analogue of the Bakry--\'Emery identity ---
consistent with the subset inversion \eqref{eq:gamma-N-inversion}
--- is the natural all-$|S|$ refinement and is left as part of the
open problem.
\item[$(\mathrm v)$] \emph{Semigroup-level adequacy.} When
$\mathcal N=\mathcal N_{\Phi}$ is the natural noise of a stochastic
flow of kernels $\Phi$ in the sense of \cite{LeJanRaimond2004}, the
two-point form $\Gamma_{\mathcal N}^{(2)}$ restricted to product
cylindrical functionals
$F_{f,g}^{(x,y)}:=f(X_t^{(x)})g(X_t^{(y)})$ recovers the two-point
carr\'e-du-champ $\Gamma^{(2)}$ of \eqref{eq:gamma-twopoint} on
product observables $f\otimes g$, and, more generally, each
$\Gamma_{\mathcal N,S}$ restricted to $|S|$-fold product cylindrical
functionals recovers the generator-level interaction defect
$\Gamma_S$ of \cref{def:N-point-cdc} on product observables
$\prod_{i\in S}f_i$.
\end{itemize}
\end{program}

\Cref{prog:open} sharpens the open problem at three levels. At the
level of the primitive datum, $(\mathrm i)$ asks for an
intrinsically-specified positive energy kernel that does not pass
through a first-chaos generator, organized into the all-$|S|$
generator-level interaction hierarchy of
\cref{sec:twopoint-hierarchy} and headed by the bilinear two-point
form $\Gamma_{\mathcal N}^{(2)}$. At the
level of the closed adjoint pair, $(\mathrm{ii})$--$(\mathrm{iv})$
ask for $(\delta_{\mathcal N},D_{\mathcal N})$ consistent with the
two-point layer via the Bakry--\'Emery identity, with the Leibniz
defect of $D_{\mathcal N}$ as the explicit algebraic remainder, and
with the higher layers $\Gamma_{\mathcal N,S}$ ($|S|\ge 3$) as the
natural all-$|S|$ refinement of the recovery criterion via the
subset inversion \eqref{eq:gamma-N-inversion}. At the level of
consistency with an underlying flow, $(\mathrm v)$ asks the
primitive hierarchy to recover the generator-level interaction
defects of \cref{sec:twopoint,sec:twopoint-hierarchy} on $|S|$-fold
product cylindrical functionals. The two pressure points
$(\mathrm A)$ and $(\mathrm B)$ of \cref{rem:diagnostic-fork} are
the two locations where the construction must depart from the
classical pattern; $(\mathrm v)$ at the two-point level provides
the first concrete adequacy criterion against which a candidate
construction can be tested.

Two outcomes are consistent with the framework of the present paper
and with the literature it references. We state both symmetrically.

\medskip\noindent
\emph{Positive outcome.} There exists a class of nontrivial black
noises $\mathcal N$ for which a tuple
$(\{\Gamma_{\mathcal N,S}\}_{|S|\ge 2},\delta_{\mathcal N},D_{\mathcal N})$
as in \cref{prog:open} can be constructed. Then the energy-first
It\^o calculus genuinely extends beyond classical innovation: the
two-point form $\Gamma_{\mathcal N}^{(2)}$ supplies the positive
energy kernel that the seed lemma rules out at the first-chaos
level, the higher layers $\Gamma_{\mathcal N,S}$, $|S|\ge 3$,
record the genuinely $|S|$-body generator-level algebraic content
that pairwise energy does not capture, $\delta$ can be primitive
over a nonclassical noise, and observable It\^o identities driven by
such noises admit a Malliavin-type lift in which the first chaos is
absent and the algebra of the calculus is carried by a nonclassical
seed (in the senses of
\cref{rem:seed-interpretation,sec:cdc-seed,sec:twopoint-hierarchy})
and, possibly, a nonzero Leibniz defect.

\medskip\noindent
\emph{Negative outcome.} No tuple as in \cref{prog:open} exists on
any nontrivial black noise. Equivalently, every attempt at a closed
primitive divergence either reduces (after closure) to an
interval-additive deterministic seed and collapses by
\cref{lem:seed}, or fails to satisfy one of
$(\mathrm i)$--$(\mathrm v)$ for other reasons. Then a nontrivial
linearizable / classical subnoise is \emph{necessary} for the
energy-first It\^o calculus, and the trilogy's framework is
intrinsically tied to noises $\mathcal N$ with
$\mathcal N^{\mathrm{cl}}\neq\{\emptyset,\Omega\}$.

\medskip
We make no conjecture as to which outcome holds. The available
evidence from Tsirelson's program on the rigidity of first chaos in
nonclassical noises~\cite{Tsirelson2004,TsirelsonVershik1998} is
consistent with, but does not prove, the negative outcome. The
purpose of \cref{prog:open} is to formulate the problem at a level
of specificity that makes either outcome accessible to a precise
attack: by exhibiting one of the five conditions to fail, or by
constructing an example in which all five hold.

\begin{remark}[On stating partial answers]\label{rem:partial-answers}
A partial answer to \cref{prog:open} need not produce the full
hierarchy or the full tuple
$(\{\Gamma_{\mathcal N,S}\}_{|S|\ge 2},\delta_{\mathcal N},D_{\mathcal N})$.
Useful partial results include: completion of
\cref{prop:defect-sticky} from the verification data of
\cref{rem:lejanraimond-todo}, making the two-point carr\'e-du-champ
$\Gamma^{(2)}_{\theta}$ of the Brownian sticky flow explicit; an
explicit construction of
$(\delta_{\mathcal N_{\Phi}},D_{\mathcal N_{\Phi}})$ on a
Le Jan--Raimond flow consistent with the two-point carr\'e-du-champ
in the sense of \cref{rem:defect-shadow}; a precise identification
of which of $(\mathrm i)$--$(\mathrm v)$ in \cref{prog:open} forces
a negative answer in a given attempted construction; a worked
example of a nonclassical seed in the senses of
\cref{rem:seed-interpretation} or
\cref{sec:cdc-seed,sec:seed-taxonomy}, even without the full
carr\'e-du-champ-divergence-derivative tuple attached; or the
identification of a nontrivial higher interaction defect
$\Gamma_{\mathcal N,S}$ for some $|S|\ge 3$ on a noise whose
two-point layer has been computed, in the sense of
\cref{def:N-point-cdc}.
\end{remark}

%==============================================================
\section{Examples}\label{sec:examples}
%==============================================================

The examples below are referenced in the body of this paper and in
\cref{rem:black-no-centered-additive}. We restate them briefly, with
attribution, and add no new proofs. The IBM paragraph is included
explicitly as a non-example: it lies on a classical noise.

\subsection{Tsirelson's sticky and coalescing flows}\label{sec:ex-sticky}

Le Jan and Raimond~\cite{LeJanRaimond2004} construct stochastic flows
of kernels --- in particular, the sticky and coalescing Brownian
flows --- whose natural noise is black. The construction proceeds by
exhibiting an explicit Tsirelson--Vershik noise structure on the
flow and showing that every additive integrable family on this noise
vanishes; we refer the reader to~\cite{LeJanRaimond2004}, especially
\S 5, for the details. These flows are arguably the cleanest known
example of black noise with an explicit tensor factorization, and
they are the natural candidates for testing partial answers to
\cref{prog:open}.

\subsection{The noise of the Brownian web}\label{sec:ex-bweb}

The Brownian web is a system of coalescing one-dimensional Brownian
motions starting from every space-time point; it was constructed by
T\'oth--Werner and Fontes--Isopi--Newman--Ravishankar, and surveyed
by Schertzer, Sun and Swart~\cite{SSS2017}. Its associated noise is
widely understood to be black in the framework of Tsirelson; see
\cite{Tsirelson2004} and subsequent work, with further exposition
and related two-dimensional black-noise constructions discussed in
\cite{SSS2017}. The one-point motion under the Brownian web is
itself a Brownian motion, but the full web noise is not generated
by any one-parameter L\'evy generator: blackness, in this sense,
forbids it. We cite this example as a non-trivial illustration that
black noises can arise from constructions whose marginal behavior
is classical.

\subsection{The noise of two-dimensional critical percolation}\label{sec:ex-percolation}

Schramm and Smirnov~\cite{SchrammSmirnov2011} prove that the scaling
limit of two-dimensional critical percolation, viewed as a random
generalized function or as a noise on the plane, is black. This is
one of the few examples of a black noise arising directly from a
critical lattice model in dimension greater than one. The result
relies on conformal invariance and on a Bernoulli-type
characterization of the scaling limit; we refer
to~\cite{SchrammSmirnov2011} for the proof and for the precise form
of the noise axioms used.

\subsection{IBM as a non-black example}\label{sec:ex-ibm}

Iterated Brownian motion $\IBM_t=X_{|Y_t|}$, with two-sided Brownian
scenery $X$ and independent Brownian clock $Y$, has natural noise
$\mathcal N_{(X,Y)}=\mathcal N_X\otimes\mathcal N_Y$ in the sense of
\eqref{eq:L2-tensor}. Both factors are classical Wiener noises with
rich first chaos, and the tensor product is classical. The
obstructions identified in the IBM trilogy
\cite{EnergyFirst,FCDC,BidegreeAtlas} are not at the level of the
noise algebra; they live, respectively, in the $L^2$ closure of the
mollified pathwise correction on a first-$X$-chaos sector
(\cite{EnergyFirst}), in the dual $H^{-1}$ closure of that sector
(\cite{FCDC}), and in the cross-sector summability of the higher
bidegree atlas (\cite{BidegreeAtlas}). \Cref{prog:open} is therefore
not addressed by, and does not subsume, any of those results.

%==============================================================
\section{Discussion and non-claims}\label{sec:discussion}
%==============================================================

We close by restating the non-claims of the paper in their final form.

\begin{enumerate}
\item[$(\mathrm 1)$] \emph{No new black-noise theorem is proved.}
The proposition of \cref{sec:proposition} restates Tsirelson's
classification of black noise in the ladder language of this paper;
the examples of \cref{sec:examples} are cited, not reproved.

\item[$(\mathrm 2)$] \emph{No primitive divergence over a black noise
is constructed.} \Cref{prog:open} is genuinely open; partial answers
of the form described in \cref{rem:partial-answers} are welcome.

\item[$(\mathrm 3)$] \emph{IBM is not a black-noise example.} Its
natural noise is the classical tensor product
$\mathcal N_X\otimes\mathcal N_Y$ of two Wiener noises, and its
trilogy obstructions are intra-classical.

\item[$(\mathrm 4)$] \emph{Black noise blocks classical innovation
reduction, not necessarily primitive divergence.} The proposition
of \cref{sec:proposition} forbids the strong reduction of
\cref{def:cir}; it does not forbid an energy lift in the sense of
\cref{def:energy-lift}. The two-level pressure-point analysis of
\cref{sec:seed,sec:leibniz} sharpens but does not resolve the
distinction.

\item[$(\mathrm 5)$] \emph{The Leibniz defect is a diagnostic, not a
proved necessity.} \Cref{prop:diagnostic} kills only the
interval-additive deterministic seed in black noise; it does not
force the adjoint derivative $D_{\mathcal N}$ to be non-first-order.
Whether blackness forces a nonzero Leibniz defect on every separating
test algebra is part of \cref{prog:open}$(\mathrm{iv})$ and is open.

\item[$(\mathrm 6)$] \emph{The hierarchy is a conceptual extension,
not a construction.} The generator-level interaction hierarchy
$\{\Gamma_{\mathcal N,S}\}_{|S|\ge 2}$ of
\cref{sec:twopoint-hierarchy} is introduced as the natural
all-$|S|$ extension of the two-point seed, at the infinitesimal
generator level rather than as a finite-time partition-cumulant
family (see \cref{rem:partition-vs-generator}). We do not claim
that the hierarchy exists, intrinsically specified from
$\mathcal N$, on any nontrivial black noise; we do not claim that it
determines $\delta_{\mathcal N}$ even when it exists; we prove no
representation theorem linking the hierarchy to a closed adjoint
pair. The two-point member $\Gamma_{\mathcal N}^{(2)}$ remains the
primary candidate seed and the first concrete layer to be
constructed in any partial answer to \cref{prog:open}.
\end{enumerate}

The intended role of the present paper in the operator-derivative
program is to identify the first two pressure points that any
extension of the framework to nonclassical noise must address:
confinement of the deterministic-coefficient interval-additive seed
to the first chaos (\cref{lem:seed}), and the Leibniz defect of the
adjoint derivative on a test algebra (\cref{def:leibniz-defect}). In
the operator-derivative framework, black noise stands beside IBM as
a first-class example along a complementary axis: IBM stresses the
framework horizontally, inside a classical noise, through failures
of $L^2$ closure on specific bidegree sectors; black noise stresses
it vertically, by removing the classical first-chaos input entirely.
The trilogy~\cite{EnergyFirst,FCDC,BidegreeAtlas} addresses the
horizontal axis; the present paper formulates the vertical axis as
an open research problem (\cref{prog:open}), with two concrete
locations at which a positive answer must depart from the classical
pattern and two concrete locations at which a negative answer would
have to be established.

The conclusion is not that black noise defeats the operator
derivative. The conclusion is that black noise isolates the point at
which the operator derivative becomes genuinely representation-free.

%==============================================================
\begin{thebibliography}{99}

\bibitem{EnergyFirst}
R.~Fontes,
\emph{Energy-First Stochastic Calculus: Iterated Brownian Motion as a Test Case}.
Preprint.

\bibitem{FCDC}
R.~Fontes,
\emph{First-Chaos Dual Closure of the IBM It\^o Correction}.
Preprint.

\bibitem{BidegreeAtlas}
R.~Fontes,
\emph{A Bidegree Atlas for the IBM It\^o Correction}.
Preprint.

\bibitem{EthierKurtz}
S.~N.~Ethier and T.~G.~Kurtz,
\emph{Markov Processes: Characterization and Convergence}.
Wiley, New York, 1986.

\bibitem{LastPenrose}
G.~Last and M.~Penrose,
\emph{Lectures on the Poisson Process}.
Cambridge University Press, Cambridge, 2017.

\bibitem{LeJanRaimond2004}
Y.~Le~Jan and O.~Raimond,
Flows, coalescence and noise.
\emph{Ann. Probab.} \textbf{32} (2004), no.~2A, 1247--1315.

\bibitem{Nualart}
D.~Nualart,
\emph{The Malliavin Calculus and Related Topics}, 2nd ed.
Springer, Berlin, 2006.

\bibitem{SchrammSmirnov2011}
O.~Schramm and S.~Smirnov,
On the scaling limits of planar percolation.
\emph{Ann. Probab.} \textbf{39} (2011), no.~5, 1768--1814.

\bibitem{SSS2017}
E.~Schertzer, R.~Sun and J.~M.~Swart,
The Brownian web, the Brownian net, and their universality.
In: \emph{Advances in Disordered Systems, Random Processes and Some Applications}, P.~Contucci and C.~Giardin\`a (eds.), Cambridge University Press, 2017.

\bibitem{Tsirelson2004}
B.~Tsirelson,
Nonclassical stochastic flows and continuous products.
\emph{Probab. Surv.} \textbf{1} (2004), 173--298.

\bibitem{TsirelsonVershik1998}
B.~S.~Tsirelson and A.~M.~Vershik,
Examples of nonlinear continuous tensor products of measure spaces and non-Fock factorizations.
\emph{Rev. Math. Phys.} \textbf{10} (1998), no.~1, 81--145.

\bibitem{BakryGentilLedoux}
D.~Bakry, I.~Gentil and M.~Ledoux,
\emph{Analysis and Geometry of Markov Diffusion Operators}.
Grundlehren der mathematischen Wissenschaften, vol.~348, Springer, 2014.

\bibitem{FontesFracBE}
R.~Fontes,
\emph{Fractional Laplacian carr\'e-du-champ and covariance kernels: a Bakry--\'Emery view}.
In preparation.

\end{thebibliography}

\end{document}
